\title[La justicia de Dios]{\bfseries{La justicia de Dios}}
\subtitle{``Justo es \SENOR{} en todos sus caminos'' (Sal 145:17)}
\author{}\institute{}\date{}
\titlegraphic{\includegraphics[width=0.80\linewidth,height=0.36\textheight,keepaspectratio]{portada.jpg}\\[-1pt]
  {\fontsize{5}{6}\selectfont\color{gris}\textit{Nota.} Elaboración propia: Sal 145:17a en texto hebreo consonántico (BHS).}}
\setbeamerfont{section in toc}{size=\small}
\begin{document}
\frame{\titlepage}

\begin{frame}\frametitle{Contenidos}\tableofcontents\end{frame}

% ------------------------------------------------------------------
\section{Contexto}

\begin{frame}{Conocer al que hace justicia}
\begin{columns}[c,onlytextwidth]
\begin{column}{0.54\textwidth}
\begin{cita}Mas alábese en esto el que se hubiere de alabar: en entenderme y conocerme, que yo soy \SENOR, que hago misericordia, juicio y justicia en la tierra; porque estas cosas quiero, dice \SENOR.\fuente{(Jer 9:24 [9:23], RVR1960)}\end{cita}
\begin{center}\heb{עֹשֶׂה חֶסֶד מִשְׁפָּט וּצְדָקָה בָּאָרֶץ}\end{center}
\end{column}
\begin{column}{0.44\textwidth}
\centering
\begin{tikzpicture}[font=\scriptsize]
 \node[caja,text width=1.45cm,align=center] (a) at (0,0) {\heb{חֶסֶד}\\ \emph{jésed}\\ misericordia};
 \node[draw=dorado,line width=1.5pt,fill=azul,text=papiro,rounded corners=3pt,text width=1.45cm,align=center,inner sep=5pt] (b) at (1.85,0) {\heb{מִשְׁפָּט}\\ \emph{mishpat}\\ juicio};
 \node[draw=dorado,line width=1.5pt,fill=azul,text=papiro,rounded corners=3pt,text width=1.45cm,align=center,inner sep=5pt] (c) at (3.7,0) {\heb{צְדָקָה}\\ \emph{tsedaqá}\\ justicia};
 \draw[decorate,decoration={brace,amplitude=6pt,mirror},gris] (1.0,-0.9)--(4.5,-0.9) node[midway,below=7pt,font=\tiny,text=gris]{el tema de este estudio};
\end{tikzpicture}\par\medskip
\small De las tres cosas que Dios ``hace'' en la tierra, \destacar{dos} son justicia. Conocer a Dios es conocer al que hace justicia.
\end{column}
\end{columns}
\end{frame}

\begin{frame}{Tres movimientos}
\centering
\begin{tikzpicture}[font=\scriptsize,
  c/.style={caja,text width=3.2cm,align=center},
  d/.style={draw=dorado,line width=1.5pt,fill=azul,text=papiro,rounded corners=3pt,text width=3.2cm,align=center,inner sep=6pt},
  f/.style={-{Latex[length=2.6mm]},dorado,line width=1.4pt}]
\node[d] (a) at (0,0) {\small\textbf{I. Dios es justo}\\[2pt] su justicia es su ser\\ ``Justo es \SENOR{} en todos sus caminos'' (Sal 145:17)};
\node[c] (b) at (4.0,0) {\small\textbf{II. Decantada en sus mandamientos}\\[2pt] \emph{mishpatim}, \emph{edot}, \emph{juqqim}\\ ``tu ley la verdad'' (Sal 119:142)};
\node[c] (c) at (8.0,0) {\small\textbf{III. Aplicada a nosotros en Cristo}\\[2pt] ``me rodeó de manto de justicia'' (Is 61:10)};
\draw[f] (a)--(b); \draw[f] (b)--(c);
\end{tikzpicture}\par\bigskip
\small Primero, Dios es justo en sí mismo. Segundo, esa justicia eterna \destacar{baja} a la historia y se pone por escrito en los mandamientos dados a Israel. Tercero, como nadie la cumplió, Dios la \destacar{cumplió en Cristo} y la puso sobre su pueblo como un manto.
\end{frame}

\begin{frame}{Un atributo de identidad}
\begin{columns}[c,onlytextwidth]
\begin{column}{0.52\textwidth}\small
\begin{tcolorbox}[colback=papiro,colframe=dorado,boxrule=1.2pt,arc=3pt]
\centering\textbf{Dios es justo en todos sus caminos}
\end{tcolorbox}
\begin{itemize}
\item Atributo de \destacar{Identidad}: responde a \emph{¿quién es Dios?}
\item \destacar{Comunicable}: Dios la posee en perfección y pide al ser humano ``hacer justicia'' (Mi 6:8).
\item Perfección eterna del \emph{ser}, no solo un \emph{acto}: Dios era justo antes de que existiera un pecador.
\end{itemize}
\end{column}
\begin{column}{0.46\textwidth}
\scriptsize\renewcommand{\arraystretch}{1.2}
\begin{tabular}{@{}>{\RaggedRight\arraybackslash}p{1.9cm}>{\RaggedRight\arraybackslash}p{3.6cm}@{}}
\toprule
\textbf{Clase} & \textbf{Atributos} \\ \midrule
Ser & Uno, Amor, Espíritu \\
\destacar{Identidad} & Santo, Eterno, Soberano, \destacar{Justo}, Inmutable, Omnipotente, Omnipresente, Omnisciente \\
Relación & Sabiduría, Verdadero, Misericordioso, Fiel, Celo, Ira santa, Bondad, Paciencia, Gracia \\
\bottomrule
\end{tabular}
\end{column}
\end{columns}\par\smallskip
\precision{clasificación según el plan de la serie y el material del autor (``Justicia (comunicable)''); la fórmula ``perfección eterna, comunicada como reflejo'' es propuesta de esta presentación.}
\end{frame}

\begin{frame}{``Justicia'' en la RVR1960}
\begin{columns}[c,onlytextwidth]
\begin{column}{0.56\textwidth}
\scriptsize\renewcommand{\arraystretch}{1.18}
\begin{tabular}{@{}>{\RaggedRight\arraybackslash}p{3.4cm}rrr@{}}
\toprule
\textbf{Palabra} & \textbf{AT} & \textbf{NT} & \textbf{Total} \\ \midrule
justicia(s), justiciera & 277 & 101 & 378 \\
justo, justa, justos, justas & 250 & 91 & 341 \\
justificar, justificación\ldots & 21 & 44 & 65 \\
justamente & 3 & 2 & 5 \\ \midrule
\textbf{Familia positiva} & \textbf{554} & \textbf{238} & \textbf{792} \\
injusticia, injusto\ldots & 22 & 33 & 55 \\
\bottomrule
\end{tabular}
\end{column}
\begin{column}{0.42\textwidth}
\begin{tikzpicture}
\foreach \y/\l/\a/\n in {0/justicia/2.77/1.01, -0.75/justo/2.50/0.91, -1.5/justificar/0.21/0.44}{
 \node[anchor=east,font=\scriptsize] at (0,\y) {\l};
 \fill[azul] (0.1,\y-0.17) rectangle ++(\a*1.1,0.34);
 \fill[dorado] (0.1+\a*1.1,\y-0.17) rectangle ++(\n*1.1,0.34);}
\node[font=\tiny,text=azul] at (1.2,-2.2) {\rule{6pt}{6pt} AT};
\node[font=\tiny,text=dorado!80!black] at (2.4,-2.2) {\rule{6pt}{6pt} NT};
\end{tikzpicture}\par\smallskip
\small El NT usa ``\destacar{justificar}'' el doble que el AT: el Dios justo es también el que justifica.
\end{column}
\end{columns}\par\smallskip
\precision{conteo propio de apariciones en la RVR1960 sin títulos de sección (792 en 717 versículos); verificar en e-Sword.}
\end{frame}

\begin{frame}[plain]
\vfill\centering
{\color{dorado}\rule{0.55\linewidth}{1.2pt}}\\[8pt]
{\large\color{gris}Primera parte}\\[4pt]
{\LARGE\bfseries\color{azul}Dios es justo}\\[8pt]
{\small\itshape\color{gris}``Justo es \SENOR{} en todos sus caminos'' (Sal 145:17)}\\[8pt]
{\color{dorado}\rule{0.55\linewidth}{1.2pt}}
\vfill
\end{frame}

% ------------------------------------------------------------------
\section{Salmo 145}

\begin{frame}{Texto ancla: Sal 145:17}
\fotolado{panel_salmo145.jpg}{Elaboración propia: Sal 145:17 en texto hebreo consonántico (BHS).}
\begin{columns}[c,onlytextwidth]
\begin{column}{0.58\textwidth}
\begin{cita}Justo es \SENOR{} en todos sus caminos, y misericordioso en todas sus obras.\fuente{(Sal 145:17, RVR1960)}\end{cita}
\begin{center}\heb{צַדִּיק יְהוָה בְּכָל־דְּרָכָיו}\\[2pt]\heb{וְחָסִיד בְּכָל־מַעֲשָׂיו}\end{center}
\small El Salmo 145 es un \destacar{acróstico}: cada verso empieza con una letra del alfabeto. El v.~17 es la línea de la letra \heb{צ} \emph{tsade}, y empieza con \emph{tsaddiq}, ``justo''.
\end{column}
\begin{column}{0.40\textwidth}\end{column}
\end{columns}
\end{frame}

\begin{frame}{Sal 145:17: la palabra}
\scriptsize\renewcommand{\arraystretch}{1.15}
\begin{tabular}{@{}>{\RaggedRight\arraybackslash}p{2.8cm}>{\RaggedRight\arraybackslash}p{1.1cm}>{\RaggedRight\arraybackslash}p{8.0cm}@{}}
\toprule
\textbf{Término} & \textbf{Raíz} & \textbf{Sentido} \\ \midrule
\heb{צַדִּיק} \emph{tsaddiq} & \heb{צדק} & ``justo'': forma intensiva. El que está conforme a la norma y la cumple. Va primero: ``justo es el \SENOR''. \\
\heb{בְּכָל־דְּרָכָיו} & \heb{דרך} & ``en todos sus caminos'': su conducta, su manera de actuar. \\
\heb{וְחָסִיד} \emph{vejasid} & \heb{חסד} & ``y leal, fiel en amor'' (de \emph{jésed}); RVR1960 ``misericordioso''; King James ``holy''. \\
\heb{בְּכָל־מַעֲשָׂיו} & \heb{עשׂה} & ``en todas sus obras''. \\
\bottomrule
\end{tabular}\par\medskip
\begin{columns}[t,onlytextwidth]
\begin{column}{0.49\textwidth}\small
\textbf{LXX:} \grc{δίκαιος κύριος ἐν πάσαις ταῖς ὁδοῖς αὐτοῦ καὶ ὅσιος ἐν πᾶσι τοῖς ἔργοις αὐτοῦ} (Sal 144:17 LXX).
\end{column}
\begin{column}{0.49\textwidth}\small
La King James traduce \emph{holy}, ``santo'', y así lo cita Pink a propósito de la \destacar{santidad} (2020, cap. 8).
\end{column}
\end{columns}
\end{frame}

\begin{frame}{Dos palabras en paralelo}
\centering
\begin{tikzpicture}[font=\small]
 \node[draw=dorado,line width=1.5pt,fill=azul,text=papiro,rounded corners=3pt,text width=3.3cm,align=center,inner sep=6pt] (a) at (0,1.2) {\textbf{Justo} \heb{צַדִּיק}};
 \node[caja,text width=3.3cm,align=center,font=\small] (b) at (5.2,1.2) {en \textbf{todos} sus caminos};
 \node[draw=dorado,line width=1.5pt,fill=azul,text=papiro,rounded corners=3pt,text width=3.3cm,align=center,inner sep=6pt] (c) at (0,-0.3) {\textbf{Misericordioso} \heb{חָסִיד}};
 \node[caja,text width=3.3cm,align=center,font=\small] (d) at (5.2,-0.3) {en \textbf{todas} sus obras};
 \draw[dorado,line width=1.2pt,{Latex[length=2.5mm]}-{Latex[length=2.5mm]}] (a)--(c);
 \draw[dorado,line width=1.2pt,{Latex[length=2.5mm]}-{Latex[length=2.5mm]}] (b)--(d);
 \draw[gris,line width=0.8pt] (a)--(b); \draw[gris,line width=0.8pt] (c)--(d);
\end{tikzpicture}\par\medskip
\begin{itemize}\small
\item La justicia de Dios \destacar{no está sola}: camina junto a su amor leal. No son fuerzas en tensión, sino dos descripciones del mismo actuar.
\item \heb{כָּל} \emph{kol}, ``todo'', aparece \textbf{17 veces} en el salmo: no hay zona de la conducta de Dios donde la justicia quede suspendida.
\end{itemize}
\precision{conteo de \emph{kol} propio, en la BHS.}
\end{frame}

\begin{frame}{El verso en su salmo}
\begin{columns}[c,onlytextwidth]
\begin{column}{0.54\textwidth}
\begin{cita}Cercano está \SENOR{} a todos los que le invocan, a todos los que le invocan de veras. Cumplirá el deseo de los que le temen; oirá asimismo el clamor de ellos, y los salvará. \SENOR{} guarda a todos los que le aman, mas destruirá a todos los impíos.\fuente{(Sal 145:18--20, RVR1960)}\end{cita}
\end{column}
\begin{column}{0.44\textwidth}\small
\begin{itemize}
\item v.~7: ``cantarán tu \destacar{justicia}'' (\heb{וְצִדְקָתְךָ יְרַנֵּנוּ}).
\item v.~8: Éx 34:6, el Nombre revelado en el Sinaí.
\item vv.~18--20: la misma justicia \destacar{salva} al que clama y \destacar{destruye} al impío.
\end{itemize}
\end{column}
\end{columns}\par\medskip
\precision{(textual) el TM no tiene verso para la letra \emph{nun}; la LXX (144:13a) y 11QPs\textsuperscript{a} sí: ``Fiel es Dios en sus palabras, y santo en todas sus obras'', con la misma segunda mitad que el v.~17. La RVR1960 sigue el TM.}
\end{frame}

\begin{frame}{``Todos sus caminos'': Moisés}
\begin{columns}[c,onlytextwidth]
\begin{column}{0.52\textwidth}
\begin{cita}El es la Roca, cuya obra es perfecta, porque todos sus caminos son rectitud; Dios de verdad, y sin ninguna iniquidad en él; es justo y recto.\fuente{(Dt 32:4, RVR1960)}\end{cita}
\begin{center}\small\heb{כִּי כָל־דְּרָכָיו מִשְׁפָּט}\\\heb{צַדִּיק וְיָשָׁר הוּא}\end{center}
\end{column}
\begin{column}{0.46\textwidth}\scriptsize
\begin{itemize}
\item \heb{כָל־דְּרָכָיו מִשְׁפָּט}: ``todos sus caminos son \emph{juicio}''.
\item \heb{וְאֵין עָוֶל}: ``no hay \emph{ável}'', torcedura del derecho.
\item \heb{צַדִּיק וְיָשָׁר}: ``justo y \emph{recto}'' (lo que no se tuerce).
\item Este verso abre la oración judía ante la muerte (\emph{Tsidduq haddín}).
\item Ap 15:3: ``justos y verdaderos son \destacar{tus caminos}'': ``el cántico de Moisés\ldots{} y el cántico del Cordero''.
\end{itemize}
\end{column}
\end{columns}
\end{frame}

% ------------------------------------------------------------------
\section{La palabra}

\begin{frame}{La raíz \emph{tsadaq}}
\begin{columns}[c,onlytextwidth]
\begin{column}{0.64\textwidth}
\scriptsize\renewcommand{\arraystretch}{1.18}
\begin{tabular}{@{}>{\RaggedRight\arraybackslash}p{1.8cm}>{\RaggedRight\arraybackslash}p{2.3cm}r>{\RaggedRight\arraybackslash}p{1.9cm}@{}}
\toprule
\textbf{Término} & \textbf{Sentido} & \textbf{Veces} & \textbf{RVR1960} \\ \midrule
\heb{צַדִּיק} \emph{tsaddiq} & justo & 205 & justo, justos \\
\heb{צְדָקָה} \emph{tsedaqá} & justicia; acto que salva & 157 & justicia, honradez \\
\heb{צֶדֶק} \emph{tsedeq} & justicia; lo exacto & 118 & justicia, justo \\
\heb{צָדַק} (verbo) & ser justo; justificar & 41 & justificar(se) \\ \midrule
\textbf{Total} & & \textbf{521} & \\
\bottomrule
\end{tabular}
\end{column}
\begin{column}{0.34\textwidth}
\begin{tikzpicture}
\foreach \y/\l/\v in {0/tsaddiq/205, -0.6/tsedaqá/157, -1.2/tsedeq/118, -1.8/verbo/41}{
 \node[anchor=east,font=\tiny\itshape] at (0,\y) {\l};
 \fill[azul] (0.1,\y-0.17) rectangle ++(\v/95,0.34);
 \node[anchor=west,font=\tiny] at (0.15+\v/95,\y) {\v};}
\end{tikzpicture}\par\smallskip
\small Sentido base: \destacar{conformidad con una norma}.
\end{column}
\end{columns}\par\smallskip
\precision{conteo propio por números de Strong (H6662--6666) en la RVR1960-Strong, cotejado con la BHS; coincide con BDB y HALOT. Verificar en e-Sword.}
\end{frame}

\begin{frame}{Balanzas justas}
\begin{columns}[c,onlytextwidth]
\begin{column}{0.52\textwidth}
\begin{cita}Balanzas justas, pesas justas y medidas justas tendréis. Yo \SENOR{} vuestro Dios, que os saqué de la tierra de Egipto.\fuente{(Lv 19:36, RVR1960)}\end{cita}
\begin{center}\small\heb{מֹאזְנֵי צֶדֶק אַבְנֵי־צֶדֶק אֵיפַת צֶדֶק וְהִין צֶדֶק}\end{center}
\scriptsize\emph{tsedeq} cuatro veces: la balanza, las pesas (``piedras''), el \emph{efá} (áridos) y el \emph{hin} (líquidos).
\end{column}
\begin{column}{0.46\textwidth}
\centering
\begin{tikzpicture}[scale=0.8]
 \draw[azuloscuro,line width=2pt] (0,0)--(0,2.6); \fill[azuloscuro] (-0.6,0) rectangle (0.6,0.15);
 \draw[azul,line width=1.8pt] (-2,2.6)--(2,2.6); \fill[dorado] (0,2.6) circle (3.5pt);
 \foreach \x in {-2,2}{\draw[gris] (\x,2.6)--(\x-0.6,1.5); \draw[gris] (\x,2.6)--(\x+0.6,1.5);
   \draw[azul,line width=1.4pt,fill=papiro] (\x-0.7,1.5) arc (180:360:0.7 and 0.3);}
 \fill[gris] (-2.25,1.45) rectangle (-1.75,1.75);
 \node[font=\tiny,text=gris] at (-2,1.95) {pesa};
 \node[font=\scriptsize\bfseries,text=azul] at (0,3.05) {\heb{מֹאזְנֵי צֶדֶק}};
\end{tikzpicture}\par\smallskip
\small Una balanza exacta \destacar{refleja el carácter de Dios} en el mercado: ``Peso y balanzas justas son de \SENOR'' (Pr 16:11).
\end{column}
\end{columns}
\end{frame}

\begin{frame}{Usos concretos de la raíz}
\scriptsize\renewcommand{\arraystretch}{1.25}
\begin{tabular}{@{}>{\RaggedRight\arraybackslash}p{2.0cm}>{\RaggedRight\arraybackslash}p{3.0cm}>{\RaggedRight\arraybackslash}p{7.0cm}@{}}
\toprule
\textbf{Texto} & \textbf{Hebreo} & \textbf{Qué muestra} \\ \midrule
Sal 23:3 & \heb{מַעְגְּלֵי־צֶדֶק} & ``sendas de justicia'': caminos rectos, que llevan a donde deben. \\
Gn 30:33 & \heb{צִדְקָתִי} & ``mi honradez'' (Jacob ante Labán): honradez comprobable en un negocio. \\
Gn 38:26 & \heb{צָדְקָה מִמֶּנִּי} & ``Más justa es ella que yo'' (Judá de Tamar): ella tiene la razón. \\
Dt 25:15 & \heb{אֶבֶן שְׁלֵמָה וָצֶדֶק} & ``pesa exacta y justa''. \\
Dt 16:20 & \heb{צֶדֶק צֶדֶק תִּרְדֹּף} & ``la justicia, la justicia seguirás'': el mandato a los jueces. \\
\bottomrule
\end{tabular}\par\bigskip
\begin{tcolorbox}[colback=papiro,colframe=dorado,boxrule=1pt,arc=3pt]
\small La justicia hebrea es \destacar{relacional}: cumplir lo que una relación exige, conforme a una norma. En último término, la norma es el carácter de Dios.
\end{tcolorbox}
\end{frame}

\begin{frame}{\emph{Mishpat}: el juicio}
\begin{columns}[c,onlytextwidth]
\begin{column}{0.50\textwidth}\small
\begin{itemize}
\item \heb{מִשְׁפָּט} \emph{mishpat}, de \heb{שׁפט} ``juzgar, gobernar'': \textbf{425} veces en la BHS.
\item Es la decisión del juez, el derecho que resulta y la conducta conforme a él.
\item El par \heb{צְדָקָה וּמִשְׁפָּט} es fórmula fija: Gn 18:19; Sal 89:14; Jer 9:24; 23:5; Is 9:7; Am 5:24.
\end{itemize}
\end{column}
\begin{column}{0.48\textwidth}
\scriptsize\renewcommand{\arraystretch}{1.12}
\begin{tabular}{@{}>{\RaggedRight\arraybackslash}p{3.0cm}r@{}}
\toprule
\textbf{RVR1960 traduce \emph{mishpat}} & \textbf{Veces} \\ \midrule
juicio(s) & $\approx$147 \\
decretos & 36 \\
derecho & 32 \\
justicia & 16 \\
ley(es), ordenanzas, costumbre, causa, rito & resto \\
\bottomrule
\end{tabular}
\end{column}
\end{columns}\par\medskip
\begin{cita}Pero corra el juicio como las aguas, y la justicia como impetuoso arroyo.\fuente{(Am 5:24, RVR1960)}\end{cita}
\precision{conteo propio (BHS: 425; etiquetado Strong H4941: 417).}
\end{frame}

\begin{frame}{Las ``justicias'' del \SENOR}
\begin{columns}[c,onlytextwidth]
\begin{column}{0.50\textwidth}\small
\heb{צִדְקוֹת יְהוָה} \emph{tsidqot YHWH}, ``las justicias del \SENOR'', son sus \destacar{actos de liberación}:
\begin{itemize}
\item ``los \textbf{triunfos} de \SENOR'' (Jue 5:11)
\item ``los \textbf{hechos de salvación} que \SENOR{} ha hecho'' (1 S 12:7)
\item ``para que conozcas las \textbf{justicias} de \SENOR'' (Mi 6:5)
\end{itemize}
\end{column}
\begin{column}{0.48\textwidth}
\begin{cita}\SENOR{} ha hecho notoria su salvación; a vista de las naciones ha descubierto su justicia.\fuente{(Sal 98:2, RVR1960)}\end{cita}
\scriptsize\heb{גִּלָּה צִדְקָתוֹ} ``reveló su justicia''; LXX \grc{ἀπεκάλυψεν τὴν δικαιοσύνην αὐτοῦ}: la frase volverá en Ro 1:17.
\end{column}
\end{columns}\par\medskip
\small Para Israel, la justicia de Dios no es solo el tribunal que castiga: es sobre todo \destacar{la acción con que cumple su pacto} (Is 46:13; 51:5).
\end{frame}

\begin{frame}{El espectro griego}
\begin{columns}[c,onlytextwidth]
\begin{column}{0.56\textwidth}
\scriptsize\renewcommand{\arraystretch}{1.1}
\begin{tabular}{@{}>{\RaggedRight\arraybackslash}p{1.8cm}>{\RaggedRight\arraybackslash}p{1.9cm}r>{\RaggedRight\arraybackslash}p{1.4cm}@{}}
\toprule
\textbf{Término} & \textbf{Sentido} & \textbf{NT} & \textbf{Texto} \\ \midrule
\grc{δικαιοσύνη} & justicia & 92 & Ro 1:17 \\
\grc{δίκαιος} & justo & 79 & Ro 3:26 \\
\grc{δικαιόω} & justificar & 39 & Ro 4:5 \\
\grc{δικαίωμα} & requisito justo & 10 & Ro 8:4 \\
\grc{δικαίωσις} & justificación & 2 & Ro 4:25 \\
\grc{δικαιοκρισία} & justo juicio & 1 & Ro 2:5 \\
\grc{ἀδικία} & injusticia & 25 & Ro 1:18 \\
\bottomrule
\end{tabular}
\end{column}
\begin{column}{0.42\textwidth}\small
\grc{δικαιοσύνη}: \textbf{58} de 92 en las cartas de Pablo (\textbf{34} en Romanos).\par\medskip
\precision{conteo propio por Strong en el NA27; sentidos según BDAG.}
\end{column}
\end{columns}
\end{frame}

\begin{frame}{Del hebreo a Pablo}
\centering
\begin{tikzpicture}[font=\small,
  c/.style={draw=azul,line width=0.8pt,fill=papiro,fill opacity=0.92,text opacity=1,rounded corners=3pt,text width=3.1cm,align=center,inner sep=5pt},
  f/.style={-{Latex[length=2.6mm]},dorado,line width=1.4pt}]
\node[c] (h) at (0,0) {\textbf{Hebreo}\\ \heb{צְדָקָה} \emph{tsedaqá}\\ \heb{צַדִּיק} \emph{tsaddiq}\\ \heb{הִצְדִּיק} \emph{hitsdiq}};
\node[c] (g) at (4.1,0) {\textbf{Septuaginta}\\ \grc{δικαιοσύνη}\\ \grc{δίκαιος}\\ \grc{δικαιόω}};
\node[c] (p) at (8.2,0) {\textbf{Pablo}\\ ``la justicia de Dios'' (Ro 1:17)\\ ``el justo, y el que justifica'' (Ro 3:26)};
\draw[f] (h)--(g); \draw[f] (g)--(p);
\node[font=\footnotesize\itshape,text=gris,text width=11cm,align=center] at (4.1,-1.75) {Pablo escribe con el vocabulario de la Biblia griega de Israel, y con él trae el sentido hebreo: una norma, un juicio y una acción que salva.};
\end{tikzpicture}
\end{frame}

% ------------------------------------------------------------------
\section{El Nombre}

\begin{frame}{``\SENOR, justicia nuestra''}
\fotolado{panel_jeremias23.jpg}{Elaboración propia: Jer 23:6b en texto hebreo consonántico (BHS).}
\begin{columns}[c,onlytextwidth]
\begin{column}{0.58\textwidth}
\begin{cita}He aquí que vienen días, dice \SENOR, en que levantaré a David renuevo justo, y reinará como Rey\ldots{} y hará juicio y justicia en la tierra. En sus días será salvo Judá\ldots{} y este será su nombre con el cual le llamarán: \SENOR, justicia nuestra.\fuente{(Jer 23:5--6, RVR1960)}\end{cita}
\begin{center}\heb{יְהוָה צִדְקֵנוּ}\end{center}
\scriptsize La justicia tiene un \destacar{nombre propio}, y es un nombre \destacar{mesiánico}.
\end{column}
\begin{column}{0.40\textwidth}\end{column}
\end{columns}
\end{frame}

\begin{frame}{Jer 23:5--6: la palabra}
\scriptsize\renewcommand{\arraystretch}{1.15}
\begin{tabular}{@{}>{\RaggedRight\arraybackslash}p{2.8cm}>{\RaggedRight\arraybackslash}p{1.1cm}>{\RaggedRight\arraybackslash}p{8.0cm}@{}}
\toprule
\textbf{Término} & \textbf{Raíz} & \textbf{Sentido} \\ \midrule
\heb{צֶמַח צַדִּיק} & \heb{צמח} & ``renuevo justo'': el brote del tronco cortado de David (Is 11:1; Zac 3:8; 6:12). \\
\heb{יְהוָה צִדְקֵנוּ} & \heb{צדק} & ``el \SENOR{} (es) nuestra justicia'': frase nominal usada como nombre propio. \\
\grc{Ἰωσεδεκ} (LXX) & --- & la LXX no traduce: translitera \emph{Iōsedek} (= \emph{Yehotsadaq}, ``el \SENOR{} es justo''). \\
\bottomrule
\end{tabular}\par\medskip
\begin{columns}[c,onlytextwidth]
\begin{column}{0.49\textwidth}
\centering
\begin{tikzpicture}[font=\scriptsize]
 \node[caja,text width=2.3cm,align=center,fill=gris!15] (a) at (0,0) {\heb{צִדְקִיָּהוּ}\\ Sedequías\\ ``mi justicia es el \SENOR''};
 \node[draw=dorado,line width=1.4pt,fill=azul,text=papiro,rounded corners=3pt,text width=2.3cm,align=center,inner sep=5pt] (b) at (3.1,0) {\heb{יְהוָה צִדְקֵנוּ}\\ el Rey que viene\\ ``el \SENOR, justicia nuestra''};
 \draw[-{Latex[length=2.5mm]},dorado,line width=1.2pt] (a)--(b);
\end{tikzpicture}
\end{column}
\begin{column}{0.49\textwidth}\scriptsize
Sedequías llevó el nombre y \destacar{no hizo justicia}; el Rey que viene tendrá las mismas palabras en orden inverso, y las cumplirá. En Jer 33:16 el nombre se da a Jerusalén.\par
\precision{el juego de nombres es lectura frecuente (Lundbom, 2004), no explícita.}
\end{column}
\end{columns}
\end{frame}

\begin{frame}{Tres llamados con el Nombre}
\begin{columns}[c,onlytextwidth]
\begin{column}{0.50\textwidth}\small
\begin{center}\heb{שְׁלֹשָׁה נִקְרְאוּ עַל שְׁמוֹ שֶׁל הַקָּדוֹשׁ בָּרוּךְ הוּא}\end{center}
``Tres fueron llamados con el nombre del Santo, bendito sea: los justos, el Mesías y Jerusalén'' (\emph{Bava Batra} 75b).
\end{column}
\begin{column}{0.48\textwidth}
\scriptsize\renewcommand{\arraystretch}{1.2}
\begin{tabular}{@{}>{\RaggedRight\arraybackslash}p{1.6cm}>{\RaggedRight\arraybackslash}p{3.6cm}@{}}
\toprule
\textbf{Quién} & \textbf{Prueba} \\ \midrule
Los justos & Is 43:7, ``todos los llamados de mi nombre'' \\
El Mesías & Jer 23:6, ``\SENOR, justicia nuestra'' \\
Jerusalén & Ez 48:35, ``\SENOR-sama'' \\
\bottomrule
\end{tabular}
\end{column}
\end{columns}\par\medskip
\precision{interpretación rabínica (R. Shemuel bar Najmaní en nombre de R. Yojanán).}
\end{frame}

\begin{frame}{\emph{Elohim}: el Nombre del juicio}
\begin{columns}[c,onlytextwidth]
\begin{column}{0.52\textwidth}\small
\emph{Shemot Rabá} 3:6 (R. Abbá bar Memel): ``Según mis obras soy llamado''.
\begin{itemize}
\item cuando \destacar{juzgo}: \emph{Elohim}
\item cuando combato a los impíos: \emph{Tsevaot}
\item cuando suspendo el castigo: \emph{El Shadday}
\item cuando tengo misericordia: el \SENOR
\end{itemize}
\end{column}
\begin{column}{0.46\textwidth}\small
Rashi a Gn 1:1: ``creó \emph{Elohim}'' porque Dios pensó crear el mundo con la \heb{מִדַּת הַדִּין} \emph{middat haddín} (medida de la justicia), y le asoció la de la misericordia.\par\medskip
Sal 7:11 [7:12]: \heb{אֱלֹהִים שׁוֹפֵט צַדִּיק}, ``Dios es juez justo''.
\end{column}
\end{columns}\par\medskip
\precision{principio de lectura rabínico, no regla gramatical: \emph{Elohim} aparece también en textos de misericordia.}
\end{frame}

\begin{frame}{``Dios justo y Salvador''}
\begin{columns}[c,onlytextwidth]
\begin{column}{0.54\textwidth}
\begin{cita}\ldots{}¿quién hizo oír esto desde el principio\ldots{} sino yo \SENOR? Y no hay más Dios que yo; Dios justo y Salvador; ningún otro fuera de mí.\fuente{(Is 45:21, RVR1960)}\end{cita}
\begin{center}\heb{אֵל־צַדִּיק וּמוֹשִׁיעַ}\end{center}
\end{column}
\begin{column}{0.44\textwidth}\small
\begin{itemize}
\item \emph{El tsaddiq umoshía}: justicia y salvación en un solo título.
\item ``En \SENOR{} será \destacar{justificada} y se gloriará toda la descendencia de Israel'' (Is 45:25).
\item Zac 9:9: ``tu rey vendrá a ti, \destacar{justo y salvador}''.
\item Pablo cita el pasaje en Ro 14:11 y Flp 2:10--11.
\end{itemize}
\end{column}
\end{columns}
\end{frame}

\begin{frame}{Los nombres de la justicia}
\scriptsize\renewcommand{\arraystretch}{1.15}
\begin{tabular}{@{}>{\RaggedRight\arraybackslash}p{3.1cm}>{\RaggedRight\arraybackslash}p{3.0cm}>{\RaggedRight\arraybackslash}p{2.0cm}>{\RaggedRight\arraybackslash}p{3.4cm}@{}}
\toprule
\textbf{Nombre o título} & \textbf{Hebreo / griego} & \textbf{Texto} & \textbf{Qué revela} \\ \midrule
El Juez de toda la tierra & \heb{שֹׁפֵט כָּל־הָאָרֶץ} & Gn 18:25 & Juzga a todos los pueblos \\
Dios juez justo & \heb{אֱלֹהִים שׁוֹפֵט צַדִּיק} & Sal 7:11 & \emph{Elohim}, Nombre del juicio \\
Dios de justicia & \heb{אֱלֹהֵי מִשְׁפָּט} & Is 30:18; Mal 2:17 & Espera para tener piedad \\
Dios justo y Salvador & \heb{אֵל־צַדִּיק וּמוֹשִׁיעַ} & Is 45:21 & Justicia que salva \\
\SENOR, justicia nuestra & \heb{יְהוָה צִדְקֵנוּ} & Jer 23:6; 33:16 & El Mesías, justicia del pueblo \\
Sol de justicia & \heb{שֶׁמֶשׁ צְדָקָה} & Mal 4:2 [3:20] & Justicia que sana \\
El Justo & \grc{ὁ δίκαιος} & Hch 3:14; 22:14 & Título de Cristo \\
Juez justo & \grc{ὁ δίκαιος κριτής} & 2 Ti 4:8 & Corona en aquel día \\
\bottomrule
\end{tabular}\par\medskip
\small Pablo ``vio al \destacar{Justo}'' en el camino a Damasco (Hch 22:14): a ese encuentro le debe su comprensión de la justicia de Dios.
\end{frame}

% ------------------------------------------------------------------
\section{Definición}

\begin{frame}{Dios es justicia}
\begin{columns}[c,onlytextwidth]
\begin{column}{0.52\textwidth}
\begin{cita}Porque \SENOR{} es justo, y ama la justicia; el hombre recto mirará su rostro.\fuente{(Sal 11:7, RVR1960)}\end{cita}
\begin{cita}Para anunciar que \SENOR{} mi fortaleza es recto, y que en él no hay injusticia.\fuente{(Sal 92:15, RVR1960)}\end{cita}
\end{column}
\begin{column}{0.46\textwidth}\small
\begin{itemize}
\item No hay una norma \destacar{por encima} de Dios con la cual se lo mida: \destacar{la norma es él}.
\item ``Dios de verdad, y sin ninguna iniquidad en él; es justo y recto'' (Dt 32:4).
\item Su nombre mesiánico no dice ``el \SENOR{} es justo'', sino ``\SENOR, \emph{justicia} nuestra'' (Jer 23:6).
\item Dios era justo antes de que existiera un pecador y antes de toda ley escrita.
\end{itemize}
\end{column}
\end{columns}
\end{frame}

\begin{frame}{Cualidad, no solo acto}
\small El material de la serie: \destacar{``justicia como cualidad moral, no como acto''}; Dios no se vuelve justo cuando juzga.\par\medskip
\centering
\begin{tikzpicture}[font=\scriptsize,
  c/.style={caja,text width=2.6cm,align=center},
  d/.style={draw=dorado,line width=1.5pt,fill=azul,text=papiro,rounded corners=3pt,text width=2.6cm,align=center,inner sep=5pt},
  f/.style={-{Latex[length=2.6mm]},dorado,line width=1.4pt}]
\node[c] (s) at (0,0) {\textbf{Santidad}\\ puro y separado del mal (Is 5:16)};
\node[d] (j) at (3.7,0) {\textbf{Justicia}\\ Dios es la norma y actúa conforme a ella (Sal 145:17)};
\node[c] (ju) at (7.6,0.95) {\textbf{Juicio}\\ justicia + soberanía (Sal 9:7--8)};
\node[c] (di) at (7.6,-0.95) {\textbf{Disciplina}\\ justicia + amor (Sal 119:75)};
\draw[f] (s)--(j); \draw[f] (j)--(ju.west); \draw[f] (j)--(di.west);
\node[font=\tiny\itshape,text=gris] at (3.7,-2.0) {perfección eterna (ser) \quad $\rightarrow$ \quad manifestación histórica (actuar)};
\end{tikzpicture}\par\medskip
\begin{cita}\SENOR{} de los ejércitos será exaltado en juicio, y el Dios Santo será santificado con justicia.\fuente{(Is 5:16, RVR1960)}\end{cita}
\end{frame}

\begin{frame}{Aspectos de la justicia}
\scriptsize\renewcommand{\arraystretch}{1.2}
\begin{tabular}{@{}>{\RaggedRight\arraybackslash}p{2.3cm}>{\RaggedRight\arraybackslash}p{5.0cm}>{\RaggedRight\arraybackslash}p{4.6cm}@{}}
\toprule
\textbf{Aspecto} & \textbf{Qué significa} & \textbf{Textos} \\ \midrule
En su ser & Dios es recto; no hay en él injusticia & Dt 32:4; Sal 92:15; Ro 9:14 \\
Legislativa & Da leyes justas, reflejo de su carácter & Dt 4:8; Sal 19:9; Ro 7:12 \\
Retributiva & Paga a cada uno según sus obras & Job 34:11; Ro 2:5--6; 2 Ts 1:6 \\
Remunerativa & Recompensa conforme a su promesa & He 6:10; 2 Ti 4:8 \\
\destacar{Defensora} & Hace justicia al oprimido & Dt 10:18; Sal 103:6; 146:7 \\
\destacar{Salvadora} & Cumple su pacto y libera & Jue 5:11; Is 46:13; Sal 98:2 \\
\bottomrule
\end{tabular}\par\medskip
\small Las cuatro primeras siguen las distinciones clásicas (Berkhof, 1938/1996: rectoral y distributiva, remunerativa y retributiva). Las dos últimas recogen el uso hebreo: si la justicia de Dios fuera solo retributiva, sería mala noticia para todos.
\end{frame}

\begin{frame}{Del ser de Dios a su ley y a su pueblo}
\begin{columns}[t,onlytextwidth]
\begin{column}{0.49\textwidth}
\begin{cita}Tu justicia es justicia eterna, y tu ley la verdad.\fuente{(Sal 119:142, RVR1960)}\end{cita}
\end{column}
\begin{column}{0.49\textwidth}
\begin{cita}\ldots{}me vistió con vestiduras de salvación, me rodeó de manto de justicia\ldots\fuente{(Is 61:10, RVR1960)}\end{cita}
\end{column}
\end{columns}
\centering
\begin{tikzpicture}[font=\scriptsize,
  c/.style={caja,text width=3.0cm,align=center},
  d/.style={draw=dorado,line width=1.5pt,fill=azul,text=papiro,rounded corners=3pt,text width=3.0cm,align=center,inner sep=5pt},
  f/.style={-{Latex[length=2.6mm]},dorado,line width=1.4pt}]
\node[d] (e) at (0,0) {\textbf{Dios es justo}\\ ``Tu justicia es justicia eterna''};
\node[c] (l) at (3.9,0) {\textbf{Dice lo que es justo}\\ ``y tu ley la verdad'': \emph{mishpatim}, \emph{edot}, \emph{juqqim}};
\node[c] (v) at (7.8,0) {\textbf{Viste de justicia a su pueblo}\\ ``manto de justicia'' (Is 61:10; 2 Co 5:21)};
\draw[f] (e)--(l); \draw[f] (l)--(v);
\end{tikzpicture}
\end{frame}

\begin{frame}{Definición}
\begin{tcolorbox}[colback=papiro,colframe=azul,boxrule=1pt,arc=3pt]
\small\textbf{La justicia de Dios} es la perfección eterna de su ser por la cual \textbf{Dios mismo es la norma} de todo lo recto (\emph{tsedeq}). Esa justicia se \textbf{decantó en sus mandamientos} a Israel, juicios, testimonios y estatutos que son ``justicia eterna'' (Sal 119:142, 144). Y esa justicia, que nadie cumplió, Dios la \textbf{cumplió en Cristo} y la puso sobre su pueblo como un manto (Is 61:10; 2 Co 5:21).
\end{tcolorbox}
\begin{columns}[t,onlytextwidth]
\begin{column}{0.49\textwidth}
\begin{cita}Justo es \SENOR{} en todos sus caminos, y misericordioso en todas sus obras.\fuente{(Sal 145:17)}\end{cita}
\end{column}
\begin{column}{0.49\textwidth}
\begin{cita}\ldots{}a fin de que él sea el justo, y el que justifica al que es de la fe de Jesús.\fuente{(Ro 3:26)}\end{cita}
\end{column}
\end{columns}
\precision{síntesis propuesta a partir de la definición de la serie y del esquema del autor.}
\end{frame}

% ------------------------------------------------------------------
\section{Teología}

\begin{frame}{El trono de Dios}
\begin{columns}[c,onlytextwidth]
\begin{column}{0.50\textwidth}
\begin{cita}Justicia y juicio son el cimiento de tu trono; misericordia y verdad van delante de tu rostro.\fuente{(Sal 89:14 [89:15], RVR1960)}\end{cita}
\begin{center}\scriptsize\heb{צֶדֶק וּמִשְׁפָּט מְכוֹן כִּסְאֶךָ}\\\heb{חֶסֶד וֶאֱמֶת יְקַדְּמוּ פָנֶיךָ}\end{center}
\end{column}
\begin{column}{0.48\textwidth}
\centering
\begin{tikzpicture}[font=\scriptsize]
 \fill[azul] (-2.2,0) rectangle (2.2,0.55); \node[text=papiro,font=\scriptsize\bfseries] at (0,0.27) {justicia y juicio};
 \node[font=\tiny,text=gris] at (0,-0.25) {\heb{מְכוֹן} cimiento};
 \draw[dorado,line width=1.4pt,fill=papiro] (-1.2,0.55) rectangle (1.2,1.6);
 \draw[dorado,line width=1.4pt,fill=papiro] (-1.2,1.6) rectangle (-0.9,2.6); \draw[dorado,line width=1.4pt,fill=papiro] (0.9,1.6) rectangle (1.2,2.6);
 \draw[dorado,line width=1.4pt,fill=papiro] (-1.2,2.6) -- (1.2,2.6) -- (1.2,2.8) -- (-1.2,2.8) -- cycle;
 \node[font=\tiny,text=azul] at (0,1.1) {trono};
 \draw[-{Latex[length=2.2mm]},gris,line width=1pt] (1.4,1.4)--(2.3,1.4) node[right,text width=1.0cm,font=\tiny]{misericordia y verdad};
\end{tikzpicture}\par\smallskip
\scriptsize La justicia es el \destacar{cimiento}; la misericordia y la verdad, los \destacar{heraldos} que salen al encuentro. Igual en Sal 97:2.
\end{column}
\end{columns}
\end{frame}

\begin{frame}{``La justicia y la paz se besaron''}
\begin{columns}[c,onlytextwidth]
\begin{column}{0.52\textwidth}
\begin{cita}La misericordia y la verdad se encontraron; la justicia y la paz se besaron. La verdad brotará de la tierra, y la justicia mirará desde los cielos.\fuente{(Sal 85:10--11 [85:11--12], RVR1960)}\end{cita}
\begin{center}\small\heb{צֶדֶק וְשָׁלוֹם נָשָׁקוּ}\end{center}
\end{column}
\begin{column}{0.46\textwidth}\small
\begin{itemize}
\item Oración del pueblo que volvió del exilio (vv. 1--7).
\item Los atributos \destacar{se encuentran y se besan}; la justicia ``mira desde los cielos''.
\item El material de la serie: en la cruz ``la Justicia fue satisfecha, la Misericordia fue derramada y la Santidad fue vindicada''.
\end{itemize}
\end{column}
\end{columns}\par\smallskip
\precision{lectura tipológica del autor; en su contexto el salmo pide la restauración tras el exilio. El Midrash lo lee como disputa de los atributos en la creación del hombre (\emph{Bereshit Rabá} 8:5).}
\end{frame}

\begin{frame}{El nudo: ``no justificaré al impío''}
\begin{columns}[c,onlytextwidth]
\begin{column}{0.50\textwidth}\small
\begin{itemize}
\item Éx 34:7: ``que perdona la iniquidad\ldots{} y que de ningún modo tendrá por inocente al malvado''.
\item Éx 23:7: ``porque yo \destacar{no justificaré al impío}'' (\heb{לֹא־אַצְדִּיק רָשָׁע}).
\item Pr 17:15: ``El que justifica al impío, y el que condena al justo, ambos son igualmente abominación a \SENOR''.
\end{itemize}
\end{column}
\begin{column}{0.48\textwidth}
\centering
\begin{tikzpicture}[font=\scriptsize]
 \node[caja,text width=3.6cm,align=center] (a) at (0,1.3) {\heb{מַצְדִּיק רָשָׁע} ``el que justifica al impío'' = \textbf{abominación} (Pr 17:15)};
 \node[draw=dorado,line width=1.5pt,fill=azul,text=papiro,rounded corners=3pt,text width=3.6cm,align=center,inner sep=5pt] (b) at (0,-0.6) {\grc{τὸν δικαιοῦντα τὸν ἀσεβῆ} ``aquel que justifica al impío'' = \textbf{Dios} (Ro 4:5)};
 \draw[-{Latex[length=2.5mm]},dorado,line width=1.2pt] (a)--(b) node[midway,right,font=\tiny,text=gris]{¿cómo?};
\end{tikzpicture}
\end{column}
\end{columns}\par\medskip
\small El AT deja la pregunta abierta. La responde Romanos 3:21--26.
\end{frame}

\begin{frame}{Ningún atributo cruza la justicia}
\begin{columns}[t,onlytextwidth]
\begin{column}{0.49\textwidth}\small
\textbf{\color{azul}La santidad} se muestra \destacar{en} la justicia: ``el Dios Santo será santificado con justicia'' (Is 5:16). Dios ``deshonraría sus perfecciones, si considerara justo y santo lo que no lo es'' (Pink, 2020, cap. 8).\par\medskip
\textbf{\color{azul}La misericordia} no la anula: el mismo Nombre que perdona ``de ningún modo tendrá por inocente al malvado'' (Éx 34:7); ``Dios no será injusto consigo mismo'' (Pink, 2020, cap. 14).
\end{column}
\begin{column}{0.49\textwidth}\small
\textbf{\color{azul}La gracia} no la burla: ``la gracia reine por la justicia'' (Ro 5:21). Quien rechaza la gracia queda a solas con la justicia (Pink, 2020, cap. 13).\par\medskip
\textbf{\color{azul}La soberanía} no la contradice: ``¿Que hay injusticia en Dios? En ninguna manera'' (Ro 9:14). ``Dios es justo y bueno, y siempre hace lo correcto'' (Pink, 2020, cap. 19).
\end{column}
\end{columns}\par\medskip
\small La justicia es el \destacar{cimiento del trono} (Sal 89:14): por eso el nudo no se resuelve rebajándola, sino cumpliéndola en la cruz.
\end{frame}


\begin{frame}{``¿Dónde está el Dios de justicia?''}
\begin{columns}[c,onlytextwidth]
\begin{column}{0.52\textwidth}
\begin{cita}Justo eres tú, oh \SENOR, para que yo dispute contigo; sin embargo, alegaré mi causa ante ti. ¿Por qué es prosperado el camino de los impíos\ldots?\fuente{(Jer 12:1, RVR1960)}\end{cita}
\begin{center}\small\heb{צַדִּיק אַתָּה יְהוָה כִּי אָרִיב אֵלֶיךָ}\end{center}
\end{column}
\begin{column}{0.46\textwidth}\small
\begin{itemize}
\item \heb{אָרִיב} \emph{ariv}: ``entablo pleito'', lenguaje de tribunal.
\item El profeta \destacar{no niega} la justicia de Dios: \destacar{parte de ella}.
\item La queja bíblica no es incredulidad: es apelar a Dios contra la apariencia de las cosas (cf. Mal 2:17; Sal 73).
\end{itemize}
\end{column}
\end{columns}
\end{frame}

\begin{frame}{Habacuc: ``el justo por su fe vivirá''}
\begin{columns}[c,onlytextwidth]
\begin{column}{0.50\textwidth}\small
\begin{itemize}
\item 1:4: ``sale torcida la justicia''.
\item 1:13: ``¿por qué\ldots{} callas cuando destruye el impío al más justo que él?''.
\item 2:3: ``aunque tardare, espéralo''.
\end{itemize}
\begin{cita}He aquí que aquel cuya alma no es recta, se enorgullece; mas el justo por su fe vivirá.\fuente{(Hab 2:4, RVR1960)}\end{cita}
\end{column}
\begin{column}{0.48\textwidth}
\scriptsize\renewcommand{\arraystretch}{1.15}
\begin{tabular}{@{}>{\RaggedRight\arraybackslash}p{1.1cm}>{\RaggedRight\arraybackslash}p{4.1cm}@{}}
\toprule
TM & \heb{וְצַדִּיק בֶּאֱמוּנָתוֹ יִחְיֶה} ``por \emph{su} fidelidad'' \\
LXX & \grc{ὁ δὲ δίκαιος ἐκ πίστεώς μου ζήσεται} ``por \emph{mi} fidelidad'' \\
Pablo & \grc{ὁ δὲ δίκαιος ἐκ πίστεως ζήσεται} (Ro 1:17; Gá 3:11) \\
1QpHab & los que cumplen la Torá y son fieles al Maestro de Justicia \\
\bottomrule
\end{tabular}
\end{column}
\end{columns}\par\smallskip
\small La respuesta no es una teoría, sino \destacar{una manera de vivir}: esperar con fidelidad.
\end{frame}

\begin{frame}{Entrar en el santuario; confesar}
\begin{columns}[t,onlytextwidth]
\begin{column}{0.48\textwidth}
\textbf{\color{azul}Salmo 73}\par\small
``Viendo la prosperidad de los impíos'' casi resbaló (v.~3), ``hasta que entrando en el santuario de Dios, comprendí el fin de ellos'' (v.~17).\par\medskip
La justicia de Dios se ve desde \destacar{el final} de la historia, no desde la mitad.
\end{column}
\begin{column}{0.50\textwidth}
\begin{cita}Tuya es, Señor, la justicia, y nuestra la confusión de rostro\ldots\fuente{(Dn 9:7, RVR1960)}\end{cita}
\begin{cita}Pero tú eres justo en todo lo que ha venido sobre nosotros; porque rectamente has hecho, mas nosotros hemos hecho lo malo.\fuente{(Neh 9:33, RVR1960)}\end{cita}
\scriptsize ``no\ldots{} confiados en nuestras justicias, sino en tus muchas misericordias'' (Dn 9:18).
\end{column}
\end{columns}
\end{frame}

\begin{frame}[plain]
\vfill\centering
{\color{dorado}\rule{0.55\linewidth}{1.2pt}}\\[8pt]
{\large\color{gris}Segunda parte}\\[4pt]
{\LARGE\bfseries\color{azul}La justicia de Dios decantada en sus mandamientos}\\[8pt]
{\small\itshape\color{gris}``Tu justicia es justicia eterna, y tu ley la verdad'' (Sal 119:142)}\\[8pt]
{\color{dorado}\rule{0.55\linewidth}{1.2pt}}
\vfill
\end{frame}

% ------------------------------------------------------------------
\section{La ley}

\begin{frame}{``Estatutos y juicios justos''}
\begin{columns}[c,onlytextwidth]
\begin{column}{0.50\textwidth}
\begin{cita}Y ¿qué nación grande hay que tenga estatutos y juicios justos como es toda esta ley que yo pongo hoy delante de vosotros?\fuente{(Dt 4:8, RVR1960)}\end{cita}
\begin{center}\small\heb{חֻקִּים וּמִשְׁפָּטִים צַדִּיקִם}\end{center}
\end{column}
\begin{column}{0.48\textwidth}\small
\begin{itemize}
\item Lo que Dios es (\emph{tsaddiq}, Sal 145:17), eso son sus leyes (\emph{tsaddiqim}).
\item La ley es el \destacar{carácter justo de Dios} traducido en mandatos.
\item ``Ha manifestado\ldots{} sus estatutos y sus juicios a Israel. No ha hecho así con ninguna otra de las naciones'' (Sal 147:19--20).
\end{itemize}
\end{column}
\end{columns}
\end{frame}

\begin{frame}{La tríada de Deuteronomio}
\begin{columns}[c,onlytextwidth]
\begin{column}{0.52\textwidth}
\begin{cita}Mañana cuando te preguntare tu hijo, diciendo: ¿Qué significan los testimonios y estatutos y decretos que \SENOR{} nuestro Dios os mandó?\fuente{(Dt 6:20, RVR1960)}\end{cita}
\begin{cita}Y tendremos justicia cuando cuidemos de poner por obra todos estos mandamientos\ldots\fuente{(Dt 6:25, RVR1960)}\end{cita}
\end{column}
\begin{column}{0.46\textwidth}
\scriptsize\renewcommand{\arraystretch}{1.15}
\begin{tabular}{@{}>{\RaggedRight\arraybackslash}p{1.7cm}>{\RaggedRight\arraybackslash}p{1.6cm}>{\RaggedRight\arraybackslash}p{1.7cm}@{}}
\toprule
\textbf{Hebreo} & \textbf{RVR1960} & \textbf{LXX} \\ \midrule
\heb{הָעֵדֹת} \emph{edot} & testimonios & \grc{μαρτύρια} \\
\heb{הַחֻקִּים} \emph{juqqim} & estatutos & \grc{δικαιώματα} \\
\heb{הַמִּשְׁפָּטִים} \emph{mishpatim} & decretos & \grc{κρίματα} \\
\bottomrule
\end{tabular}\par\smallskip
\small La respuesta empieza por la \destacar{redención} (``nos sacó de Egipto'', 6:21) y termina en la \destacar{justicia} (\heb{וּצְדָקָה תִּהְיֶה־לָּנוּ}, 6:25). Igual en Dt 4:45.
\end{column}
\end{columns}
\end{frame}

\begin{frame}{\emph{Mishpatim}: la vida con el prójimo}
\begin{columns}[c,onlytextwidth]
\begin{column}{0.50\textwidth}
\begin{cita}Estas son las leyes que les propondrás.\fuente{(Éx 21:1, RVR1960)}\end{cita}
\begin{center}\heb{וְאֵלֶּה הַמִּשְׁפָּטִים אֲשֶׁר תָּשִׂים לִפְנֵיהֶם}\end{center}
\begin{cita}No pervertirás el derecho de tu mendigo en su pleito\ldots{} Y no angustiarás al extranjero\ldots\fuente{(Éx 23:6, 9, RVR1960)}\end{cita}
\end{column}
\begin{column}{0.48\textwidth}\small
\begin{itemize}
\item \heb{מִשְׁפָּט}, de \heb{שׁפט}: la \destacar{decisión del juez} y el derecho que nace de ella.
\item Éx 21--23: siervos, homicidio, lesiones, daños, robo, préstamos, juicio justo, extranjero.
\item Leyes que la conciencia reconoce como justas: no matar, no robar, no torcer el juicio.
\item ``los juicios de \SENOR{} son verdad, \destacar{todos justos}'' (Sal 19:9).
\end{itemize}
\end{column}
\end{columns}
\end{frame}

\begin{frame}{La justicia que defiende al débil}
\begin{columns}[c,onlytextwidth]
\begin{column}{0.54\textwidth}
\begin{cita}Porque \SENOR{} vuestro Dios es Dios de dioses\ldots{} que no hace acepción de personas, ni toma cohecho; que hace justicia al huérfano y a la viuda; que ama también al extranjero dándole pan y vestido.\fuente{(Dt 10:17--18, RVR1960)}\end{cita}
\end{column}
\begin{column}{0.44\textwidth}\scriptsize
\begin{itemize}
\item \heb{לֹא־יִשָּׂא פָנִים}: ``no levanta rostros'', no favorece al importante.
\item \heb{עֹשֶׂה מִשְׁפַּט יָתוֹם וְאַלְמָנָה}: ``hace el derecho'' del huérfano y la viuda.
\item Sal 103:6; 146:7: ``hace justicia a los agraviados''.
\item Sal 82: Dios acusa a los jueces: ``Defended al débil y al huérfano''.
\end{itemize}
\end{column}
\end{columns}
\end{frame}

\begin{frame}{``¿No es esto conocerme a mí?''}
\begin{columns}[c,onlytextwidth]
\begin{column}{0.52\textwidth}
\begin{cita}El juzgó la causa del afligido y del menesteroso, y entonces estuvo bien. ¿No es esto conocerme a mí? dice \SENOR.\fuente{(Jer 22:16, RVR1960)}\end{cita}
\begin{center}\small\heb{הֲלוֹא־הִיא הַדַּעַת אֹתִי}\end{center}
\end{column}
\begin{column}{0.46\textwidth}
\centering
\begin{tikzpicture}[font=\scriptsize]
 \node[draw=dorado,line width=1.5pt,fill=azul,text=papiro,rounded corners=3pt,text width=3.4cm,align=center,inner sep=5pt] (a) at (0,1.2) {Conocer al Dios que \textbf{hace} justicia (Jer 9:24)};
 \node[caja,text width=3.4cm,align=center] (b) at (0,-0.4) {\textbf{hacer} justicia al pobre (Jer 22:16)};
 \draw[{Latex[length=2.5mm]}-{Latex[length=2.5mm]},dorado,line width=1.2pt] (a)--(b);
\end{tikzpicture}
\end{column}
\end{columns}\par\medskip
\begin{cita}Oh hombre, él te ha declarado lo que es bueno, y qué pide \SENOR{} de ti: solamente hacer justicia, y amar misericordia, y humillarte ante tu Dios.\fuente{(Mi 6:8, RVR1960)}\end{cita}
\end{frame}

\begin{frame}{\emph{Edot}: el testimonio}
\begin{columns}[c,onlytextwidth]
\begin{column}{0.50\textwidth}
\begin{cita}Y dio a Moisés\ldots{} dos tablas del testimonio, tablas de piedra escritas con el dedo de Dios.\fuente{(Éx 31:18, RVR1960)}\end{cita}
\begin{center}\heb{שְׁנֵי לֻחֹת הָעֵדֻת}\end{center}
\end{column}
\begin{column}{0.48\textwidth}\small
\begin{itemize}
\item \heb{עֵדוּת}, de \heb{עוד} ``atestiguar'': lo que \destacar{declara} cómo son las cosas.
\item Las tablas son ``el testimonio''; el arca, ``el arca del testimonio'', y desde allí Dios habla (Éx 25:21--22).
\item El libro de la ley junto al arca ``por testigo contra ti'' (Dt 31:26).
\item Recuerdan lo que Dios hizo: Pascua, sábado, fiestas.
\item ``Justicia eterna son tus testimonios'' (Sal 119:144).
\end{itemize}
\end{column}
\end{columns}
\end{frame}

\begin{frame}{\emph{Juqqim}: lo grabado por Dios}
\begin{columns}[c,onlytextwidth]
\begin{column}{0.50\textwidth}
\begin{cita}Mis ordenanzas pondréis por obra, y mis estatutos guardaréis, andando en ellos. Yo \SENOR{} vuestro Dios.\fuente{(Lv 18:4, RVR1960)}\end{cita}
\begin{center}\small\heb{אֶת־מִשְׁפָּטַי תַּעֲשׂוּ וְאֶת־חֻקֹּתַי תִּשְׁמְרוּ}\end{center}
\scriptsize Los \emph{mishpatim} se \emph{hacen}; los \emph{juqqim} se \emph{custodian}. ``\ldots{}los cuales haciendo el hombre, vivirá en ellos'' (18:5 = Ro 10:5; Gá 3:12).
\end{column}
\begin{column}{0.48\textwidth}\scriptsize
\begin{itemize}
\item \heb{חֹק}, de \heb{חקק} ``grabar'': ``en las palmas de las manos te tengo \destacar{esculpida}'' (Is 49:16).
\item También el orden de la creación: ``las \destacar{leyes} de la luna y de las estrellas'' (\heb{חֻקֹּת}, Jer 31:35); ``puso al mar su estatuto'' (Pr 8:29).
\item Normas cuya razón no se explica: Pascua, ``estatuto perpetuo'' (Éx 12:14); la vaca alazana, ``ordenanza de la ley'' (Nm 19:2); mezcla de hilos (Lv 19:19); ``pacto de sal'' (Nm 18:19).
\item Su razón última: ``Yo \SENOR{} vuestro Dios''.
\end{itemize}
\end{column}
\end{columns}
\end{frame}

\begin{frame}{Una ley, tres miradas}
\centering
\begin{tikzpicture}[font=\scriptsize,
  c/.style={caja,text width=3.0cm,align=center},
  d/.style={draw=dorado,line width=1.5pt,fill=azul,text=papiro,rounded corners=3pt,text width=3.0cm,align=center,inner sep=5pt},
  f/.style={-{Latex[length=2.6mm]},dorado,line width=1.4pt}]
\node[d] (j) at (3.9,1.9) {\textbf{Dios es justo}\\ \emph{tsaddiq} (Sal 145:17)};
\node[c] (m) at (0,0) {\heb{מִשְׁפָּטִים} \textbf{juicios}\\ el juez que decide\\ la vida con el prójimo\\ (Éx 21:1; Sal 19:9)};
\node[c] (e) at (3.9,0) {\heb{עֵדֹת} \textbf{testimonios}\\ el testigo que declara\\ la memoria del pacto\\ (Éx 31:18; Sal 119:144)};
\node[c] (q) at (7.8,0) {\heb{חֻקִּים} \textbf{estatutos}\\ lo grabado en piedra\\ lo fijado por Dios\\ (Lv 18:4; Nm 19:2)};
\draw[f] (j)--(m.north); \draw[f] (j)--(e.north); \draw[f] (j)--(q.north);
\end{tikzpicture}\par\medskip
\small La Biblia usa los tres términos \destacar{juntos} y a veces los intercambia (Dt 4:45; 5:1; 1 R 2:3; Neh 10:29): son tres miradas de la misma ley justa.
\end{frame}

\begin{frame}{Antes del Sinaí}
\begin{columns}[c,onlytextwidth]
\begin{column}{0.50\textwidth}
\begin{cita}\ldots{}Allí les dio estatutos y ordenanzas, y allí los probó.\fuente{(Éx 15:25, RVR1960)}\end{cita}
\begin{center}\heb{שָׁם שָׂם לוֹ חֹק וּמִשְׁפָּט}\end{center}
\begin{cita}\ldots{}y declaro las ordenanzas de Dios y sus leyes.\fuente{(Éx 18:16, RVR1960)}\end{cita}
\end{column}
\begin{column}{0.48\textwidth}\small
\begin{itemize}
\item En Mara, antes del Sinaí, Dios ya da \emph{joq} y \emph{mishpat}.
\item Moisés enseña \emph{juqqim} y \emph{torot} antes de recibir las tablas.
\item Y antes aún, Abraham ``guardó mi precepto, mis mandamientos, mis estatutos y mis leyes'' (Gn 26:5).
\item La justicia de Dios tenía forma de mandamiento antes de quedar \destacar{escrita en piedra}.
\end{itemize}
\end{column}
\end{columns}
\end{frame}

\begin{frame}{Antecedente histórico: el judaísmo y la ley}
\begin{tcolorbox}[colback=dorado!10,colframe=dorado!85!azul,boxrule=0.6pt,arc=2pt,fontupper=\scriptsize,title={\scriptsize No es la justicia de Dios declarada en la Biblia: es cómo la tradición posterior organizó la ley},coltitle=azuloscuro,colbacktitle=dorado!30]
\begin{itemize}\setlength{\itemsep}{1pt}
\item \textbf{Halajá} (\heb{הֲלָכָה}, de \heb{הלך} ``andar''): ``el camino por el cual se anda''; luego, todo el sistema legal judío.
\item \textbf{Torá oral}: Misná (Rabí Yehudá haNasí, c. 200 e.\,c.) + Guemará ($\approx$300 años) = Talmud.
\item \textbf{613 mandamientos}: 365 negativos + 248 positivos (R. Simlai, \emph{Makot} 23b). Los profetas los ``resumieron'': David en 11 (Sal 15), Isaías en 6, Miqueas en 3 (Mi 6:8), Habacuc en 1: ``el justo por su fe vivirá'' (\emph{Makot} 24a).
\item \textbf{Tres clases} (desde Dt 6:20): \emph{mishpatim} (juicios: lo que la razón reconoce: no matar, no robar), \emph{edot} (testimonios: sábado, fiestas), \emph{juqqim} (estatutos sin razón explicada: vaca roja, mezcla de telas). \emph{Yomá} 67b: los \emph{mishpatim} ``si no estuvieran escritos, sería justo escribirlos''.
\end{itemize}
\end{tcolorbox}
\precision{clasificaciones rabínicas y medievales; la Biblia usa los términos juntos, sin un sistema fijo (Sal 19; 119; 1 R 2:3). Cotejar \emph{Makot} 23b--24a y \emph{Yomá} 67b en Sefaria.}
\end{frame}

% ------------------------------------------------------------------
\section{Sal 19 y 119}

\begin{frame}{``La ley de \SENOR{} es perfecta''}
\begin{columns}[c,onlytextwidth]
\begin{column}{0.50\textwidth}
\begin{cita}La ley de \SENOR{} es perfecta, que convierte el alma; el testimonio de \SENOR{} es fiel, que hace sabio al sencillo. Los mandamientos de \SENOR{} son rectos, que alegran el corazón; el precepto de \SENOR{} es puro, que alumbra los ojos. El temor de \SENOR{} es limpio, que permanece para siempre; los juicios de \SENOR{} son verdad, todos justos.\fuente{(Sal 19:7--9 [19:8--10], RVR1960)}\end{cita}
\end{column}
\begin{column}{0.48\textwidth}\small
\begin{itemize}
\item vv. 1--6: Dios en la \destacar{creación} (los cielos, el sol).
\item vv. 7--11: Dios en su \destacar{ley}: seis líneas iguales: nombre + ``de \SENOR'' + cualidad + efecto.
\item La serie termina en \heb{צָדְקוּ יַחְדָּו}, ``son justos todos juntos''.
\item ``Deseables\ldots{} más que el oro\ldots{} dulces más que miel'' (v. 10).
\end{itemize}
\end{column}
\end{columns}
\end{frame}

\begin{frame}{Sal 19: los seis nombres de la ley}
\scriptsize\renewcommand{\arraystretch}{1.18}
\begin{tabular}{@{}>{\RaggedRight\arraybackslash}p{2.6cm}>{\RaggedRight\arraybackslash}p{1.8cm}>{\RaggedRight\arraybackslash}p{2.9cm}>{\RaggedRight\arraybackslash}p{4.6cm}@{}}
\toprule
\textbf{Nombre} & \textbf{RVR1960} & \textbf{Cualidad} & \textbf{Efecto} \\ \midrule
\heb{תּוֹרָה} \emph{torá} & ley & \heb{תְּמִימָה} perfecta & \heb{מְשִׁיבַת נָפֶשׁ} ``hace volver el alma'' \\
\heb{עֵדוּת} \emph{edut} & testimonio & \heb{נֶאֱמָנָה} fiel & \heb{מַחְכִּימַת פֶּתִי} hace sabio al simple \\
\heb{פִּקּוּדִים} \emph{piqqudim} & mandamientos & \heb{יְשָׁרִים} rectos & \heb{מְשַׂמְּחֵי־לֵב} alegran el corazón \\
\heb{מִצְוָה} \emph{mitsvá} & precepto & \heb{בָּרָה} pura & \heb{מְאִירַת עֵינָיִם} alumbra los ojos \\
\heb{יִרְאָה} \emph{yir'á} & temor & \heb{טְהוֹרָה} limpio & \heb{עוֹמֶדֶת לָעַד} permanece para siempre \\
\heb{מִשְׁפָּטִים} \emph{mishpatim} & juicios & \heb{אֱמֶת} verdad & \heb{צָדְקוּ יַחְדָּו} \destacar{todos justos} \\
\bottomrule
\end{tabular}\par\medskip
\begin{columns}[t,onlytextwidth]
\begin{column}{0.49\textwidth}\small
\heb{תְּמִימָה} \emph{temimá}: la raíz de ``la Roca, cuya obra es \emph{perfecta}'' (Dt 32:4). La ley tiene la perfección de su autor.
\end{column}
\begin{column}{0.49\textwidth}\small
Cada cualidad de la ley es una cualidad de Dios: la ley es un \destacar{retrato moral} de su autor.
\end{column}
\end{columns}
\end{frame}

\begin{frame}{El vocabulario de la ley}
\scriptsize\renewcommand{\arraystretch}{1.12}
\begin{tabular}{@{}>{\RaggedRight\arraybackslash}p{2.5cm}>{\RaggedRight\arraybackslash}p{3.9cm}>{\RaggedRight\arraybackslash}p{2.8cm}r r@{}}
\toprule
\textbf{Término} & \textbf{Raíz} & \textbf{RVR1960} & \textbf{AT} & \textbf{Sal 119} \\ \midrule
\heb{תּוֹרָה} \emph{torá} & \heb{ירה} señalar, enseñar & ley & $\approx$220 & 25 \\
\heb{מִצְוָה} \emph{mitsvá} & \heb{צוה} mandar & mandamiento, precepto & $\approx$180 & 22 \\
\heb{חֹק} / \heb{חֻקָּה} & \heb{חקק} grabar, fijar & estatuto & $\approx$230 & 22 \\
\heb{מִשְׁפָּט} \emph{mishpat} & \heb{שׁפט} juzgar & juicio, decreto & $\approx$420 & 23 \\
\heb{עֵדוּת} / \heb{עֵדָה} & \heb{עוד} atestiguar & testimonio & $\approx$85 & 23 \\
\heb{פִּקּוּדִים} & \heb{פקד} encargar & mandamientos & 24 & 21 \\
\heb{מִשְׁמֶרֶת} & \heb{שׁמר} guardar & precepto & $\approx$78 & --- \\
\heb{דָּבָר} / \heb{אִמְרָה} & palabra, dicho & palabra, dichos & --- & 24 / 19 \\
\bottomrule
\end{tabular}\par\medskip
\precision{conteo propio por Strong en la RVR1960-Strong (aproximado en el AT; el del Sal 119 coincide con los comentarios). Verificar en e-Sword.}
\end{frame}


\begin{frame}{El Salmo 119 y la tríada}
\scriptsize\renewcommand{\arraystretch}{1.2}
\begin{tabular}{@{}>{\RaggedRight\arraybackslash}p{1.7cm}r>{\RaggedRight\arraybackslash}p{6.0cm}>{\RaggedRight\arraybackslash}p{3.0cm}@{}}
\toprule
\textbf{Término} & \textbf{Veces} & \textbf{Lo que dice el salmo} & \textbf{Textos} \\ \midrule
\emph{Mishpatim} & 23 & ``tus justos juicios'' (``los juicios de tu justicia''); ``tus juicios son justos'' & vv. 7, 62, 75, 106, 160, 164 \\
\emph{Edot} & 23 & ``Justicia eterna son tus testimonios''; ``mis delicias y mis consejeros'' & vv. 24, 138, 144 \\
\emph{Juqqim} & 22 & ``Enséñame tus estatutos'' (estribillo); ``cánticos fueron para mí tus estatutos'' & vv. 12, 26, 54, 64, 68, 124, 135, 171 \\
\emph{Mitsvot} & 22 & ``todos tus mandamientos son justicia'' & v. 172 \\
\emph{Torá} & 25 & ``Tu justicia es justicia eterna, y tu ley la verdad''; ``¡cuánto amo yo tu ley!'' & vv. 142, 97 \\
\bottomrule
\end{tabular}\par\medskip
\begin{columns}[t,onlytextwidth]
\begin{column}{0.49\textwidth}\small
22 estrofas de 8 versos, una por letra. Casi todos los versos nombran la ley.
\end{column}
\begin{column}{0.49\textwidth}\small
\heb{לַמְּדֵנִי חֻקֶּיךָ}, ``enséñame tus estatutos'' (v. 12): no pide las razones, pide que Dios se los \destacar{enseñe}.
\end{column}
\end{columns}
\precision{conteo propio por Strong; verificar en e-Sword.}
\end{frame}

\begin{frame}{La estrofa \emph{tsade} del Sal 119}
\begin{columns}[c,onlytextwidth]
\begin{column}{0.46\textwidth}
\begin{cita}Justo eres tú, oh \SENOR, y rectos tus juicios. Tus testimonios, que has recomendado, son rectos y muy fieles.\fuente{(Sal 119:137--138, RVR1960)}\end{cita}
\begin{center}\small\heb{צַדִּיק אַתָּה יְהוָה וְיָשָׁר מִשְׁפָּטֶיךָ}\end{center}
\scriptsize Acróstico: 22 estrofas de ocho versos. La estrofa \heb{צ} (vv. 137--144) está dedicada a la justicia, como la línea \heb{צ} del Sal 145.
\end{column}
\begin{column}{0.52\textwidth}
\scriptsize\renewcommand{\arraystretch}{1.15}
\begin{tabular}{@{}>{\RaggedRight\arraybackslash}p{0.7cm}>{\RaggedRight\arraybackslash}p{4.9cm}@{}}
\toprule
\textbf{v.} & \textbf{RVR1960} \\ \midrule
137 & ``Justo eres tú\ldots{} y rectos tus juicios'' \\
138 & ``mandaste \emph{tsedeq} tus testimonios'' (hebreo) \\
142 & ``Tu justicia es justicia eterna, y tu ley la verdad'' \\
144 & ``Justicia eterna son tus testimonios'' \\
160 & ``eterno es todo juicio de tu justicia'' \\
172 & ``todos tus mandamientos son justicia'' \\
\bottomrule
\end{tabular}\par\smallskip
``Siete veces al día te alabo a causa de tus \destacar{justos juicios}'' (v. 164).
\end{column}
\end{columns}
\end{frame}

\begin{frame}{1 Reyes 2:3: andar en sus caminos}
\begin{columns}[c,onlytextwidth]
\begin{column}{0.52\textwidth}
\begin{cita}Guarda los preceptos de \SENOR{} tu Dios, andando en sus caminos, y observando sus estatutos y mandamientos, sus decretos y sus testimonios, de la manera que está escrito en la ley de Moisés\ldots\fuente{(1 R 2:3, RVR1960)}\end{cita}
\end{column}
\begin{column}{0.46\textwidth}
\scriptsize\renewcommand{\arraystretch}{1.12}
\begin{tabular}{@{}>{\RaggedRight\arraybackslash}p{2.2cm}>{\RaggedRight\arraybackslash}p{3.0cm}@{}}
\toprule
\heb{מִשְׁמֶרֶת יְהוָה} & lo que hay que custodiar \\
\heb{לָלֶכֶת בִּדְרָכָיו} & ``andar en sus caminos'' (Sal 145:17) \\
\heb{חֻקֹּתָיו} & estatutos (\emph{juqqim}) \\
\heb{מִצְוֹתָיו} & mandamientos (\emph{mitsvot}) \\
\heb{מִשְׁפָּטָיו} & decretos (\emph{mishpatim}) \\
\heb{עֵדְוֹתָיו} & testimonios (\emph{edot}) \\
\bottomrule
\end{tabular}\par\smallskip
\small La justicia de Dios decantada en la ley se vuelve \destacar{el camino} del rey y del pueblo.
\end{column}
\end{columns}
\end{frame}

% ------------------------------------------------------------------
\section{En Israel}

\begin{frame}{Abraham: ``haciendo justicia y juicio''}
\begin{columns}[c,onlytextwidth]
\begin{column}{0.54\textwidth}
\begin{cita}Porque yo sé que mandará a sus hijos y a su casa después de sí, que guarden el camino de \SENOR, haciendo justicia y juicio\ldots\fuente{(Gn 18:19, RVR1960)}\end{cita}
\begin{center}\small\heb{לַעֲשׂוֹת צְדָקָה וּמִשְׁפָּט}\end{center}
\end{column}
\begin{column}{0.44\textwidth}\small
\begin{itemize}
\item Primera aparición del par \emph{tsedaqá umishpat}: describe la \destacar{misión de la familia de Abraham}.
\item ``El camino del \SENOR'' es hacer justicia: el lenguaje del Sal 145:17.
\item Dios revela a Abraham lo que hará con Sodoma para que aprenda \destacar{cómo juzga Dios}.
\end{itemize}
\end{column}
\end{columns}
\end{frame}

\begin{frame}{``El Juez de toda la tierra''}
\fotolado{panel_genesis18.jpg}{Elaboración propia: Gn 18:25b en texto hebreo consonántico (BHS).}
\begin{columns}[c,onlytextwidth]
\begin{column}{0.58\textwidth}
\begin{cita}Y se acercó Abraham y dijo: ¿Destruirás también al justo con el impío?\fuente{(Gn 18:23, RVR1960)}\end{cita}
\begin{cita}Lejos de ti el hacer tal, que hagas morir al justo con el impío, y que sea el justo tratado como el impío; nunca tal hagas. El Juez de toda la tierra, ¿no ha de hacer lo que es justo?\fuente{(Gn 18:25, RVR1960)}\end{cita}
\end{column}
\begin{column}{0.40\textwidth}\end{column}
\end{columns}
\end{frame}

\begin{frame}{Gn 18:23--25: la palabra}
\scriptsize\renewcommand{\arraystretch}{1.15}
\begin{tabular}{@{}>{\RaggedRight\arraybackslash}p{2.8cm}>{\RaggedRight\arraybackslash}p{1.1cm}>{\RaggedRight\arraybackslash}p{8.0cm}@{}}
\toprule
\textbf{Término} & \textbf{Raíz} & \textbf{Sentido} \\ \midrule
\heb{וַיִּגַּשׁ} & \heb{נגשׁ} & ``y se acercó'': como quien se presenta ante un juez o ante el altar. \\
\heb{הַאַף תִּסְפֶּה} & \heb{ספה} & ``¿acaso barrerás?'': destruir de una vez, sin distinguir. \\
\heb{חָלִלָה לְּךָ} & \heb{חלל} & ``¡lejos de ti!'', lit. ``¡profanación para ti!'': indigno de Dios (dos veces). \\
\heb{הֲשֹׁפֵט כָּל־הָאָרֶץ} & \heb{שׁפט} & ``¿el Juez de toda la tierra\ldots?'': primer uso de \emph{shofet} para Dios. \\
\heb{לֹא יַעֲשֶׂה מִשְׁפָּט} & \heb{שׁפט} & ``¿no hará juicio?'': \emph{shofet} y \emph{mishpat}, la misma raíz. \\
\bottomrule
\end{tabular}\par\medskip
\small La pregunta de Abraham \destacar{no duda} de la justicia de Dios: \destacar{apela} a ella. Tratar igual al justo y al impío sería ``profanación''.
\end{frame}

\begin{frame}{De cincuenta a diez}
\centering
\begin{tikzpicture}[font=\small]
\draw[line width=1.6pt,azul,-{Latex[length=3mm]}] (0,0)--(11.6,0);
\foreach \x/\n/\v in {0.5/50/18:26, 2.4/45/18:28, 4.3/40/18:29, 6.2/30/18:30, 8.1/20/18:31, 10.0/10/18:32}{
 \fill[dorado] (\x,0) circle (5pt);
 \node[above=6pt,font=\small\bfseries,text=azul] at (\x,0) {\n};
 \node[below=6pt,font=\scriptsize,text=gris] at (\x,0) {Gn \v};}
\draw[decorate,decoration={brace,amplitude=6pt,mirror},gris] (0.3,-0.75)--(10.2,-0.75) node[midway,below=7pt,font=\footnotesize\itshape,text=gris] {``No la destruiré, si hallare allí\ldots{}'' (seis respuestas)};
\end{tikzpicture}\par\medskip
\begin{columns}[t,onlytextwidth]
\begin{column}{0.49\textwidth}\small
\begin{itemize}
\item Abraham insiste con humildad: ``soy polvo y ceniza'' (v.~27).
\item Empieza en la \destacar{justicia} (no destruir al justo) y termina en la \destacar{misericordia} (perdonar a todos por unos pocos, v.~26).
\end{itemize}
\end{column}
\begin{column}{0.49\textwidth}\small
\begin{itemize}
\item No había diez justos; pero ``Dios se acordó de Abraham, y envió fuera a Lot'' (Gn 19:29; 2 P 2:7).
\item El Juez juzgó a Sodoma \destacar{y} libró al justo.
\end{itemize}
\end{column}
\end{columns}\par\smallskip
\precision{``de la justicia a la misericordia'' es lectura propia; la relación v.~19 / v.~25 es frecuente en los comentarios (Wenham, 1994).}
\end{frame}

\begin{frame}{``Guardó mis mandamientos''}
\begin{columns}[c,onlytextwidth]
\begin{column}{0.50\textwidth}
\begin{cita}por cuanto oyó Abraham mi voz, y guardó mi precepto, mis mandamientos, mis estatutos y mis leyes.\fuente{(Gn 26:5, RVR1960)}\end{cita}
\begin{center}\small\heb{וַיִּשְׁמֹר מִשְׁמַרְתִּי מִצְוֹתַי חֻקּוֹתַי וְתוֹרֹתָי}\end{center}
\end{column}
\begin{column}{0.48\textwidth}
\scriptsize\renewcommand{\arraystretch}{1.12}
\begin{tabular}{@{}>{\RaggedRight\arraybackslash}p{1.8cm}>{\RaggedRight\arraybackslash}p{3.5cm}@{}}
\toprule
\heb{מִשְׁמַרְתִּי} & ``mi custodia'': RVR ``mi precepto'' \\
\heb{מִצְוֹתַי} & ``mis mandamientos'' \\
\heb{חֻקּוֹתַי} & ``mis estatutos'' (lo fijado) \\
\heb{תוֹרֹתָי} & ``mis instrucciones'', ``mis leyes'' \\
\bottomrule
\end{tabular}\par\smallskip
\scriptsize Primera vez que se juntan los nombres de la ley, y \destacar{antes del Sinaí}. \emph{Mitsvot}: ``Vete de tu tierra'' (12:1), ``toma ahora tu hijo'' (22:2). \emph{Juqqot}: la circuncisión, ``pacto perpetuo'' (17:13). \emph{Mishpat}: ``haciendo justicia y juicio'' (18:19).\par\precision{ejemplos de \emph{mitsvot} y \emph{juqqot} en Abraham: lectura propia.}
\end{column}
\end{columns}
\end{frame}

\begin{frame}{La justicia de Abraham}
\centering
\begin{tikzpicture}
 \foreach \x/\n/\t/\d/\r in {0/1/{Creyó}/{``le fue contado por justicia''}/{Gn 15:6}, 2.5/2/{Misión}/{``haciendo justicia y juicio''}/{Gn 18:19}, 5.0/3/{Apela}/{``el Juez de toda la tierra''}/{Gn 18:25}, 7.5/4/{Obedece}/{``obedeciste a mi voz''}/{Gn 22:18}, 10.0/5/{Guarda}/{``mis mandamientos\ldots{} y mis leyes''}/{Gn 26:5}}{
  \node[caja,anchor=north] (b\n) at (\x,0) {\parbox[t][2.1cm][t]{1.95cm}{\RaggedRight
    {\small\textbf{\textcolor{azul}{\t}}}\\[4pt]\scriptsize \d\par\vfill{\tiny\color{gris}(\r)}}};
  \node[circle,fill=azul,text=papiro,font=\scriptsize\bfseries,inner sep=1.6pt] at (b\n.north west) {\n};}
\end{tikzpicture}\par\medskip
\begin{columns}[t,onlytextwidth]
\begin{column}{0.49\textwidth}\small
\textbf{\color{azul}Pablo}: la fe ``le fue contada por justicia'' estando incircunciso (Ro 4:10); la promesa no fue ``por la ley\ldots{} sino por la justicia de la fe'' (Ro 4:13; Gá 3:17).
\end{column}
\begin{column}{0.49\textwidth}\small
\textbf{\color{azul}Santiago}: ``la fe actuó juntamente con sus obras\ldots{} y le fue contado por justicia'' (Stg 2:22--23).
\end{column}
\end{columns}\par\smallskip
\precision{la lectura rabínica de que Abraham cumplió toda la Torá (\emph{Yomá} 28b; Rashi a Gn 26:5) es antecedente histórico; aquí se sigue el sentido literal.}
\end{frame}

% ------------------------------------------------------------------
\section{Judaísmo}

\begin{frame}{El salmo de cada día}
\begin{columns}[c,onlytextwidth]
\begin{column}{0.52\textwidth}\small
\begin{itemize}
\item El Sal 145 (\emph{Ashré}) se reza \destacar{tres veces al día}: dos en la mañana, una en la tarde.
\item ``Todo el que dice \emph{Tehillá leDavid} cada día tres veces tiene asegurado que es hijo del mundo venidero'' (\emph{Berajot} 4b).
\item Razón: el orden alfabético y ``Abres tu mano\ldots'' (v.~16).
\end{itemize}
\end{column}
\begin{column}{0.46\textwidth}\small
\textbf{\color{azul}Amidá, bendición 11:}
\begin{center}\heb{מֶלֶךְ אוֹהֵב צְדָקָה וּמִשְׁפָּט}\end{center}
``Rey que ama la justicia y el derecho'' (cf. Sal 33:5). En los diez días de Rosh Hashaná a Yom Kipur: \heb{הַמֶּלֶךְ הַמִּשְׁפָּט}, ``el \destacar{Rey del juicio}'' (\emph{Berajot} 12b).
\end{column}
\end{columns}\par\medskip
\precision{fuentes judías: interpretación litúrgica y talmúdica.}
\end{frame}

\begin{frame}{``Si quieres un mundo\ldots''}
\begin{center}\heb{אִם עוֹלָם אַתָּה מְבַקֵּשׁ אֵין דִּין}\\[2pt]\heb{וְאִם דִּין אַתָּה מְבַקֵּשׁ אֵין עוֹלָם}\end{center}
\small R. Leví sobre Gn 18:25: ``Si quieres un mundo, no hay juicio (estricto); y si quieres juicio, no hay mundo. ¿Por qué tomas la cuerda por los dos extremos?'' (\emph{Bereshit Rabá} 39:6; cf. 49:9).\par\medskip
\centering
\begin{tikzpicture}[font=\scriptsize]
 \node[caja,text width=3.0cm,align=center] (a) at (0,0) {\heb{מִדַּת הַדִּין}\\ medida de la justicia};
 \node[caja,text width=3.0cm,align=center] (b) at (6,0) {\heb{מִדַּת הָרַחֲמִים}\\ medida de la misericordia};
 \draw[dorado,line width=2pt] (a)--(b) node[midway,above,font=\tiny,text=gris]{la cuerda};
 \node[font=\tiny\itshape,text=gris,text width=8cm,align=center] at (3,-0.9) {Rashi a Gn 1:1: Dios ``asoció'' la misericordia a la justicia para que el mundo subsistiera.};
\end{tikzpicture}\par
\precision{cotejar el texto completo de \emph{Bereshit Rabá} 39:6 y 49:9 en Sefaria.}
\end{frame}

\begin{frame}{\emph{Tsidduq haddín}: justificar el juicio}
\begin{columns}[c,onlytextwidth]
\begin{column}{0.52\textwidth}\small
Rabí Janiná ben Teradión, su esposa y su hija, condenados por Roma:
\begin{center}\heb{צִדְּקוּ עֲלֵיהֶם אֶת הַדִּין}\end{center}
``justificaron el juicio sobre sí'' (\emph{Avodá Zará} 18a).
\end{column}
\begin{column}{0.46\textwidth}
\scriptsize\renewcommand{\arraystretch}{1.15}
\begin{tabular}{@{}>{\RaggedRight\arraybackslash}p{1.2cm}>{\RaggedRight\arraybackslash}p{3.8cm}@{}}
\toprule
Él & ``Él es la Roca, cuya obra es perfecta'' (Dt 32:4a) \\
Esposa & ``Dios de verdad, y sin ninguna iniquidad en él'' (Dt 32:4b) \\
Hija & ``grande en consejo\ldots{} tus ojos están abiertos sobre todos los caminos'' (Jer 32:19) \\
\bottomrule
\end{tabular}
\end{column}
\end{columns}\par\medskip
\small Ante una mala noticia: \heb{בָּרוּךְ דַּיַּן הָאֱמֶת}, ``Bendito el \destacar{Juez de la verdad}'' (Misná \emph{Berajot} 9:2). La oración fúnebre \emph{Tsidduq haddín} empieza con Dt 32:4.\par
\precision{el uso fúnebre es litúrgico posterior; cotejar Misná Berajot 9:2 y 9:5 en Sefaria.}
\end{frame}

\begin{frame}{El juez y la \emph{tsedaqá}}
\begin{columns}[t,onlytextwidth]
\begin{column}{0.49\textwidth}\small
\textbf{\color{azul}Socio de Dios}\par
``Todo juez que juzga un juicio de verdad conforme a su verdad, aunque sea una sola hora\ldots{} como si se hubiera hecho \destacar{socio del Santo}, bendito sea, en la obra de la creación'' (\emph{Shabat} 10a).\par\smallskip
Dt 16:20 \emph{tsedeq tsedeq}: juicio y acuerdo justo; buscar el mejor tribunal (\emph{Sanedrín} 32b).
\end{column}
\begin{column}{0.49\textwidth}\small
\textbf{\color{azul}\emph{Tsedaqá} = caridad}\par
Dn 4:27 [4:24]: ``tus pecados redime con justicia'' (\heb{בְּצִדְקָה}). Pr 10:2 ``la justicia libra de muerte'' (\emph{Shabat} 156b).\par\smallskip
Maimónides: \emph{tsedaqá} ``denota el acto de dar a cada uno lo que le corresponde'' (\emph{Guía} III, 53). Dar al pobre \destacar{es justicia}, no generosidad opcional.
\end{column}
\end{columns}
\end{frame}

\begin{frame}{Qumrán}
\begin{columns}[c,onlytextwidth]
\begin{column}{0.56\textwidth}
\begin{cita}En cuanto a mí, si tropiezo, las misericordias de Dios serán mi salvación eterna. Si caigo por el pecado de la carne, mi justificación estará en la justicia de Dios, que permanece para siempre\ldots{} Me juzgará en la justicia de su verdad\ldots\fuente{(1QS XI, 12--14)}\end{cita}
\end{column}
\begin{column}{0.42\textwidth}\small
\begin{itemize}
\item El fundador: \heb{מוֹרֵה הַצֶּדֶק}, ``el \destacar{Maestro de Justicia}'' (CD I, 11).
\item Antes de Pablo, un judío confiesa que su justificación descansa en \heb{צִדְקַת אֵל}, ``la justicia de Dios''.
\item 4QMMT C 31: ``te será contado por justicia'' (cf. Gn 15:6).
\end{itemize}
\end{column}
\end{columns}\par\smallskip
\precision{1QS según Vermes (2004), versión propia; cotejar con García Martínez (1992) y la Leon Levy DSS Library. En Qumrán sostiene al que guarda la Torá según la comunidad; en Pablo se recibe ``aparte de la ley'' (Ro 3:21).}
\end{frame}

\begin{frame}[plain]
\vfill\centering
{\color{dorado}\rule{0.55\linewidth}{1.2pt}}\\[8pt]
{\large\color{gris}Tercera parte}\\[4pt]
{\LARGE\bfseries\color{azul}La justicia de Dios aplicada a nosotros en Cristo}\\[8pt]
{\small\itshape\color{gris}``Me rodeó de manto de justicia'' (Is 61:10)}\\[8pt]
{\color{dorado}\rule{0.55\linewidth}{1.2pt}}
\vfill
\end{frame}

% ------------------------------------------------------------------
\section{El manto}

\begin{frame}{Trapos y manto}
\begin{columns}[t,onlytextwidth]
\begin{column}{0.49\textwidth}
\begin{cita}Si bien todos nosotros somos como suciedad, y todas nuestras justicias como trapo de inmundicia\ldots\fuente{(Is 64:6, RVR1960)}\end{cita}
\small La ley muestra la justicia de Dios, y por eso muestra que \destacar{nadie la alcanza}: ``no se justificará delante de ti ningún ser humano'' (Sal 143:2).
\end{column}
\begin{column}{0.49\textwidth}
\begin{cita}En gran manera me gozaré en \SENOR, mi alma se alegrará en mi Dios; porque me vistió con vestiduras de salvación, me rodeó de manto de justicia, como a novio me atavió, y como a novia adornada con sus joyas.\fuente{(Is 61:10, RVR1960)}\end{cita}
\end{column}
\end{columns}
\begin{center}\heb{כִּי הִלְבִּישַׁנִי בִּגְדֵי־יֶשַׁע מְעִיל צְדָקָה יְעָטָנִי}\end{center}
\end{frame}

\begin{frame}{Is 61:10: la palabra}
\scriptsize\renewcommand{\arraystretch}{1.15}
\begin{tabular}{@{}>{\RaggedRight\arraybackslash}p{2.8cm}>{\RaggedRight\arraybackslash}p{1.1cm}>{\RaggedRight\arraybackslash}p{8.0cm}@{}}
\toprule
\textbf{Término} & \textbf{Raíz} & \textbf{Sentido} \\ \midrule
\heb{הִלְבִּישַׁנִי} \emph{hilbishani} & \heb{לבשׁ} & ``me \destacar{hizo vestir}'' (hifil): no se viste a sí mismo; Dios lo viste. \\
\heb{בִּגְדֵי־יֶשַׁע} & \heb{ישׁע} & ``vestiduras de salvación''. \\
\heb{מְעִיל צְדָקָה} \emph{me'il tsedaqá} & \heb{צדק} & ``manto de justicia'': la túnica larga del sumo sacerdote (Éx 28:31) y de reyes y nobles (1 S 18:4). \\
\heb{יְעָטָנִי} \emph{ye'atani} & \heb{עטה} & ``me envolvió'': la justicia lo cubre por completo. \\
\heb{כֶּחָתָן} / \heb{וְכַכַּלָּה} & --- & ``como novio\ldots{} como novia'': la ropa de bodas. \\
\bottomrule
\end{tabular}\par\medskip
\small El verbo es causativo y \destacar{Dios es el sujeto}. ``Salvación'' y ``justicia'' en paralelo: una justicia que salva (Sal 98:2; Is 46:13).\par
\precision{(textual) la LXX dice ``túnica de alegría'' (\grc{χιτῶνα εὐφροσύνης}); la RVR1960 sigue el TM, \emph{me'il tsedaqá}.}
\end{frame}

\begin{frame}{De dónde viene el manto}
\begin{columns}[c,onlytextwidth]
\begin{column}{0.52\textwidth}
\begin{cita}Y vio que no había hombre, y se maravilló que no hubiera quien se interpusiese; y lo salvó su brazo, y le afirmó su misma justicia. Pues de justicia se vistió como de una coraza, con yelmo de salvación en su cabeza\ldots\fuente{(Is 59:16--17, RVR1960)}\end{cita}
\end{column}
\begin{column}{0.46\textwidth}
\centering
\begin{tikzpicture}[font=\scriptsize]
 \node[draw=dorado,line width=1.5pt,fill=azul,text=papiro,rounded corners=3pt,text width=3.4cm,align=center,inner sep=5pt] (a) at (0,1.3) {Dios \textbf{se viste} de justicia (Is 59:17)};
 \node[caja,text width=3.4cm,align=center] (b) at (0,-0.3) {Dios \textbf{viste} de justicia a su pueblo (Is 61:10)};
 \draw[-{Latex[length=2.5mm]},dorado,line width=1.2pt] (a)--(b);
\end{tikzpicture}\par\medskip
\small El manto es \destacar{la justicia del propio Dios}. Is 61:1, ``el Espíritu\ldots{} está sobre mí'', lo leyó Jesús en Nazaret: ``Hoy se ha cumplido'' (Lc 4:18--21).
\end{column}
\end{columns}
\end{frame}

\begin{frame}{Zacarías 3: quitar las ropas viles}
\begin{columns}[c,onlytextwidth]
\begin{column}{0.52\textwidth}
\begin{cita}Y Josué estaba vestido de vestiduras viles, y estaba delante del ángel. Y habló el ángel\ldots{} diciendo: Quitadle esas vestiduras viles. Y a él le dijo: Mira que he quitado de ti tu pecado, y te he hecho vestir de ropas de gala.\fuente{(Zac 3:3--4, RVR1960)}\end{cita}
\begin{center}\small\heb{הָסִירוּ הַבְּגָדִים הַצֹּאִים מֵעָלָיו}\end{center}
\end{column}
\begin{column}{0.46\textwidth}\small
\begin{itemize}
\item El sumo sacerdote, con Satanás a su derecha para acusarlo (3:1).
\item Josué \destacar{no se limpia} ni contesta: Dios quita sus ropas, que son su pecado.
\item Dios lo viste de gala: el mismo movimiento de Is 61:10.
\item Respuesta a Is 64:6: se quitan los ``trapos'', se pone el manto.
\end{itemize}
\end{column}
\end{columns}
\end{frame}

\begin{frame}{El vestido de la justicia en la Escritura}
\scriptsize\renewcommand{\arraystretch}{1.12}
\begin{tabular}{@{}>{\RaggedRight\arraybackslash}p{1.6cm}>{\RaggedRight\arraybackslash}p{6.2cm}>{\RaggedRight\arraybackslash}p{4.2cm}@{}}
\toprule
\textbf{Texto} & \textbf{RVR1960} & \textbf{Qué muestra} \\ \midrule
Gn 3:21 & ``hizo al hombre y a su mujer túnicas de pieles, y los vistió'' & Dios cubre al culpable \\
Is 59:17 & ``de justicia se vistió como de una coraza'' & Dios se viste para salvar \\
Is 61:10 & ``me rodeó de manto de justicia'' & Dios viste a su pueblo \\
Is 64:6 & ``todas nuestras justicias como trapo de inmundicia'' & la justicia propia no cubre \\
Zac 3:4 & ``te he hecho vestir de ropas de gala'' & pecado quitado, vestido dado \\
Lc 15:22 & ``Sacad el mejor vestido, y vestidle'' & el hijo recibido por gracia \\
Mt 22:12 & ``¿cómo entraste aquí, sin estar vestido de boda?'' & sin vestido no hay boda \\
Gá 3:27 & ``de Cristo estáis revestidos'' & \destacar{el manto es Cristo} \\
Ap 19:8 & ``el lino fino es las acciones justas de los santos'' & la justicia recibida da fruto \\
\bottomrule
\end{tabular}\par
\precision{el hilo del vestido es lectura bíblico-teológica propia; Gn 3:21 como figura de la justificación es tipológico.}
\end{frame}

\begin{frame}{Revestidos de Cristo}
\begin{columns}[c,onlytextwidth]
\begin{column}{0.50\textwidth}
\begin{cita}\ldots{}porque todos los que habéis sido bautizados en Cristo, de Cristo estáis revestidos.\fuente{(Gá 3:27, RVR1960)}\end{cita}
\begin{cita}Y a ella se le ha concedido que se vista de lino fino, limpio y resplandeciente; porque el lino fino es las acciones justas de los santos.\fuente{(Ap 19:8, RVR1960)}\end{cita}
\end{column}
\begin{column}{0.48\textwidth}\small
\begin{itemize}
\item \grc{Χριστὸν ἐνεδύσασθε}: ``os vestisteis de Cristo''.
\item ``vestíos del Señor Jesucristo'' (Ro 13:14); ``vestíos del nuevo hombre, creado según Dios en la justicia'' (Ef 4:24).
\item Ap 19:8: \grc{τὰ δικαιώματα τῶν ἁγίων}, la palabra con que la LXX tradujo los \emph{juqqim} (Dt 4:45; 6:20).
\item Lo que la ley pedía, la justicia de Dios lo \destacar{produce} en los que visten a Cristo (Ro 8:4).
\end{itemize}
\end{column}
\end{columns}
\end{frame}

% ------------------------------------------------------------------
\section{En Cristo}

\begin{frame}{``No teniendo mi propia justicia''}
\begin{columns}[c,onlytextwidth]
\begin{column}{0.54\textwidth}
\begin{cita}\ldots{}y ser hallado en él, no teniendo mi propia justicia, que es por la ley, sino la que es por la fe de Cristo, la justicia que es de Dios por la fe.\fuente{(Flp 3:9, RVR1960)}\end{cita}
\end{column}
\begin{column}{0.44\textwidth}\small
\begin{itemize}
\item Antes: ``en cuanto a la justicia que es en la ley, \destacar{irreprensible}'' (Flp 3:6).
\item \grc{ἐμὴν δικαιοσύνην}: ``mi justicia''.
\item \grc{τὴν ἐκ θεοῦ δικαιοσύνην}: ``la justicia \destacar{que viene de} Dios''.
\item El cambio: el encuentro con ``el Justo'' (Hch 22:14).
\end{itemize}
\end{column}
\end{columns}
\end{frame}

\begin{frame}{La ley justa, leída con Pablo}
\begin{columns}[c,onlytextwidth]
\begin{column}{0.50\textwidth}
\begin{cita}De manera que la ley a la verdad es santa, y el mandamiento santo, justo y bueno.\fuente{(Ro 7:12, RVR1960)}\end{cita}
\begin{cita}\ldots{}para que la justicia de la ley se cumpliese en nosotros, que no andamos conforme a la carne, sino conforme al Espíritu.\fuente{(Ro 8:4, RVR1960)}\end{cita}
\end{column}
\begin{column}{0.48\textwidth}\small
\begin{itemize}
\item La ley es \grc{δικαία}, ``justa'': el problema es el ser humano (Ro 7:14; 8:3).
\item ``por la fe\ldots{} \destacar{confirmamos} la ley'' (Ro 3:31).
\item \grc{τὸ δικαίωμα τοῦ νόμου} se cumple por el Espíritu: Jer 31:33; Ez 36:27.
\item ``el cumplimiento de la ley es el amor'' (Ro 13:10).
\item ``No he venido para abrogar, sino para cumplir'' (Mt 5:17).
\end{itemize}
\end{column}
\end{columns}
\end{frame}

\begin{frame}{``La justicia de Dios se revela''}
\begin{cita}Porque no me avergüenzo del evangelio, porque es poder de Dios para salvación a todo aquel que cree; al judío primeramente, y también al griego. Porque en el evangelio la justicia de Dios se revela por fe y para fe, como está escrito: Mas el justo por la fe vivirá.\fuente{(Ro 1:16--17, RVR1960)}\end{cita}
\scriptsize\renewcommand{\arraystretch}{1.12}
\begin{tabular}{@{}>{\RaggedRight\arraybackslash}p{1.3cm}>{\RaggedRight\arraybackslash}p{3.1cm}>{\RaggedRight\arraybackslash}p{3.5cm}>{\RaggedRight\arraybackslash}p{3.4cm}@{}}
\toprule
\textbf{Texto} & \textbf{TM} & \textbf{LXX} & \textbf{Pablo} \\ \midrule
Sal 98:2 & \heb{גִּלָּה צִדְקָתוֹ} & \grc{ἀπεκάλυψεν τὴν δικαιοσύνην αὐτοῦ} & \grc{δικαιοσύνη θεοῦ\ldots{} ἀποκαλύπτεται} (Ro 1:17) \\
Sal 143:2 & \heb{לֹא־יִצְדַּק לְפָנֶיךָ כָל־חָי} & \grc{οὐ δικαιωθήσεται\ldots{} πᾶς ζῶν} & \grc{οὐ δικαιωθήσεται πᾶσα σάρξ} (Ro 3:20) \\
Gn 15:6 & \heb{וַיַּחְשְׁבֶהָ לּוֹ צְדָקָה} & \grc{ἐλογίσθη αὐτῷ εἰς δικαιοσύνην} & igual (Ro 4:3; Gá 3:6) \\
\bottomrule
\end{tabular}\par
\precision{el eco del Sal 98:2 en Ro 1:17 lo señala Hays (1989).}
\end{frame}

\begin{frame}{¿Qué es la ``justicia de Dios''?}
\small Una de las discusiones centrales de la teología del NT:\par\medskip
\scriptsize\renewcommand{\arraystretch}{1.25}
\begin{tabular}{@{}>{\RaggedRight\arraybackslash}p{3.0cm}>{\RaggedRight\arraybackslash}p{5.6cm}>{\RaggedRight\arraybackslash}p{3.4cm}@{}}
\toprule
\textbf{Lectura} & \textbf{Qué entiende} & \textbf{Apoyo} \\ \midrule
Reforma & la justicia que Dios \destacar{da} al creyente (origen) & Flp 3:9 \\
Käsemann (1969) & el \destacar{poder salvador} de Dios que reclama al mundo & Ro 1:16--17 \\
Wright (2013) & la \destacar{fidelidad de Dios a su pacto} con Israel & Ro 3:3--5; 15:8 \\
Moo (1996) & a la vez \destacar{actividad} de Dios y \destacar{don} que resulta & Ro 1:17; 3:21--22 \\
\bottomrule
\end{tabular}\par\medskip
\small La lectura hebrea de \emph{tsedaqá} (norma, juicio y acción que salva) explica por qué las tres lecturas encuentran apoyo.
\end{frame}

\begin{frame}{``El justo, y el que justifica''}
\fotolado{panel_romanos3.jpg}{Elaboración propia: Ro 3:26b en escritura uncial (texto NA28).}
\begin{columns}[c,onlytextwidth]
\begin{column}{0.58\textwidth}
\begin{cita}\ldots{}a quien Dios puso como propiciación por medio de la fe en su sangre, para manifestar su justicia, a causa de haber pasado por alto, en su paciencia, los pecados pasados, con la mira de manifestar en este tiempo su justicia, a fin de que él sea el justo, y el que justifica al que es de la fe de Jesús.\fuente{(Ro 3:25--26, RVR1960)}\end{cita}
\end{column}
\begin{column}{0.40\textwidth}\end{column}
\end{columns}
\end{frame}

\begin{frame}{Ro 3:25--26: la palabra}
\scriptsize\renewcommand{\arraystretch}{1.15}
\begin{tabular}{@{}>{\RaggedRight\arraybackslash}p{3.4cm}>{\RaggedRight\arraybackslash}p{1.4cm}>{\RaggedRight\arraybackslash}p{7.2cm}@{}}
\toprule
\textbf{Término} & \textbf{Raíz} & \textbf{Sentido} \\ \midrule
\grc{ἱλαστήριον} & \grc{ἱλάσκομαι} & ``propiciación'': en la LXX, la tapa del arca (\emph{kapporet}) rociada en Yom Kipur (Lv 16:14). \\
\grc{εἰς ἔνδειξιν τῆς δικαιοσύνης αὐτοῦ} & \grc{ἐνδείκνυμι} & ``para demostración de su justicia'': dos veces (vv. 25, 26). \\
\grc{διὰ τὴν πάρεσιν} & \grc{παρίημι} & ``por haber dejado pasar'' los pecados anteriores. \\
\grc{δίκαιον καὶ δικαιοῦντα} & \grc{δίκη} & ``justo \destacar{y} justificador'': misma raíz, unidas por ``y''. \\
\bottomrule
\end{tabular}\par\medskip
\begin{columns}[c,onlytextwidth]
\begin{column}{0.49\textwidth}\small
Dios había ``pasado por alto'' los pecados: eso ponía en duda su justicia (Pr 17:15). En la cruz trató el pecado con toda seriedad, y por eso puede justificar \destacar{sin dejar de ser justo}.
\end{column}
\begin{column}{0.49\textwidth}\small
Is 45:21: \emph{El tsaddiq umoshía}, ``Dios justo y Salvador''.\par
\precision{la conexión Ro 3:26 / Is 45:21 es lectura de varios comentaristas, no cita explícita.}
\end{column}
\end{columns}
\end{frame}

\begin{frame}{El intercambio}
\begin{columns}[c,onlytextwidth]
\begin{column}{0.50\textwidth}
\begin{cita}Al que no conoció pecado, por nosotros lo hizo pecado, para que nosotros fuésemos hechos justicia de Dios en él.\fuente{(2 Co 5:21, RVR1960)}\end{cita}
\begin{cita}\ldots{}por su conocimiento justificará mi siervo justo a muchos, y llevará las iniquidades de ellos.\fuente{(Is 53:11, RVR1960)}\end{cita}
\end{column}
\begin{column}{0.48\textwidth}
\centering
\begin{tikzpicture}[font=\scriptsize]
 \node[draw=dorado,line width=1.5pt,fill=azul,text=papiro,rounded corners=3pt,text width=2.2cm,align=center,inner sep=5pt] (c) at (0,0) {\textbf{Cristo}\\ ``el Justo''};
 \node[caja,text width=2.2cm,align=center] (n) at (3.4,0) {\textbf{nosotros}\\ ``los injustos''};
 \draw[-{Latex[length=2.5mm]},dorado,line width=1.3pt] (c.north east) to[bend left=35] node[above,font=\tiny]{justicia} (n.north west);
 \draw[-{Latex[length=2.5mm]},gris,line width=1.3pt] (n.south west) to[bend left=35] node[below,font=\tiny]{pecado} (c.south east);
\end{tikzpicture}\par\medskip
\small ``el justo por los injustos, para llevarnos a Dios'' (1 P 3:18).\par\smallskip
\scriptsize\heb{יַצְדִּיק צַדִּיק עַבְדִּי לָרַבִּים} (Is 53:11)
\end{column}
\end{columns}
\end{frame}

\begin{frame}{Jeremías en Corinto}
\begin{columns}[c,onlytextwidth]
\begin{column}{0.52\textwidth}
\begin{cita}Mas por él estáis vosotros en Cristo Jesús, el cual nos ha sido hecho por Dios sabiduría, justificación, santificación y redención; para que, como está escrito: El que se gloría, gloríese en el Señor.\fuente{(1 Co 1:30--31, RVR1960)}\end{cita}
\end{column}
\begin{column}{0.46\textwidth}
\centering
\begin{tikzpicture}[font=\scriptsize]
 \node[caja,text width=3.6cm,align=center] (a) at (0,1.5) {\textbf{Jer 9:24}: gloriarse en conocer al que hace \emph{justicia}};
 \node[caja,text width=3.6cm,align=center] (b) at (0,0.2) {\textbf{Jer 23:6}: ``\SENOR, \emph{justicia nuestra}''};
 \node[draw=dorado,line width=1.5pt,fill=azul,text=papiro,rounded corners=3pt,text width=3.6cm,align=center,inner sep=5pt] (c) at (0,-1.3) {\textbf{1 Co 1:30--31}: Cristo ``nos ha sido hecho\ldots{} justificación''};
 \draw[-{Latex[length=2.5mm]},dorado,line width=1.2pt] (a.east) to[bend left=40] (c.east);
 \draw[-{Latex[length=2.5mm]},dorado,line width=1.2pt] (b)--(c);
\end{tikzpicture}
\end{column}
\end{columns}\par\smallskip
\precision{la cita de Jer 9:24 es explícita (``como está escrito''); el vínculo con Jer 23:6 es lectura propia.}
\end{frame}

\begin{frame}{El justo juicio de Dios}
\small
\begin{itemize}
\item ``el día de la ira y de la revelación del \destacar{justo juicio} de Dios (\grc{δικαιοκρισίας}), el cual pagará a cada uno conforme a sus obras'' (Ro 2:5--6).
\item ``no hay acepción de personas para con Dios'' (Ro 2:11 = Dt 10:17).
\item ``juzgará al mundo con justicia, por aquel varón a quien designó'' (Hch 17:31): \grc{κρίνειν τὴν οἰκουμένην ἐν δικαιοσύνῃ} = Sal 9:8; 96:13; 98:9 LXX.
\item ``¿Que hay injusticia en Dios? En ninguna manera'' (Ro 9:14).
\item ``la corona de justicia, la cual me dará el Señor, juez justo, en aquel día'' (2 Ti 4:8).
\end{itemize}
\begin{tcolorbox}[colback=papiro,colframe=dorado,boxrule=1pt,arc=3pt]
\small El Juez de toda la tierra (Gn 18:25) \destacar{juzgará por medio de Cristo}.
\end{tcolorbox}
\precision{2 Timoteo: autoría discutida (cf. Ro 2:5--11; 2 Co 5:10). Hch 17:31: discurso de Pablo según Lucas.}
\end{frame}

\begin{frame}{La justicia de Dios, cumplida en Cristo}
\centering
\begin{tikzpicture}
 \foreach \x/\n/\t/\d/\r in {0/1/{Revelada}/{``la justicia de Dios se revela'' en el evangelio}/{Ro 1:17}, 2.5/2/{Demostrada}/{``el justo, y el que justifica''}/{Ro 3:26}, 5.0/3/{Dada}/{``hechos justicia de Dios en él''}/{2 Co 5:21}, 7.5/4/{Encarnada}/{``nos ha sido hecho\ldots{} justificación''}/{1 Co 1:30}, 10.0/5/{Ejercida}/{``juzgará al mundo con justicia''}/{Hch 17:31}}{
  \node[caja,anchor=north] (b\n) at (\x,0) {\parbox[t][2.4cm][t]{1.95cm}{\RaggedRight
    {\small\textbf{\textcolor{azul}{\t}}}\\[4pt]\scriptsize \d\par\vfill{\tiny\color{gris}(\r)}}};
  \node[circle,fill=azul,text=papiro,font=\scriptsize\bfseries,inner sep=1.6pt] at (b\n.north west) {\n};}
\end{tikzpicture}\par\bigskip
\small En Cristo ``habita corporalmente toda la plenitud de la Deidad'' (Col 2:9). La justicia y la paz se besaron: ``Justificados, pues, por la fe, tenemos \destacar{paz} para con Dios'' (Ro 5:1).
\end{frame}

% ------------------------------------------------------------------
\section{Vivirlo}

\begin{frame}{Vivir ante el Dios justo (1)}
\begin{columns}[t,onlytextwidth]
\begin{column}{0.49\textwidth}
\textbf{\color{azul}1. Confiar en el Juez}\par\small
``El Juez de toda la tierra, ¿no ha de hacer lo que es justo?'' (Gn 18:25). ``No os venguéis vosotros mismos\ldots{} Mía es la venganza, yo pagaré'' (Ro 12:19). ``¿Y acaso Dios no hará justicia a sus escogidos?'' (Lc 18:7).\par\medskip
\textbf{\color{azul}2. Llevarle la queja}\par\small
``Justo eres tú\ldots{} sin embargo, alegaré mi causa ante ti'' (Jer 12:1). Una apelación, no una acusación (Job 40:8).
\end{column}
\begin{column}{0.49\textwidth}
\textbf{\color{azul}3. Reconocerla en la disciplina}\par\small
``Tuya es, Señor, la justicia'' (Dn 9:7). ``Conozco, oh \SENOR, que tus juicios son justos, y que conforme a tu fidelidad me afligiste'' (Sal 119:75). ``En lugar de murmurar, le agradecemos por ello'' (Pink, 2020, cap. 10).
\end{column}
\end{columns}
\end{frame}

\begin{frame}{Vivir ante el Dios justo (2)}
\begin{columns}[t,onlytextwidth]
\begin{column}{0.49\textwidth}
\textbf{\color{azul}4. Vestirse del manto recibido}\par\small
``todas nuestras justicias como trapo de inmundicia'' (Is 64:6); ``me rodeó de manto de justicia'' (Is 61:10); ``vestíos del Señor Jesucristo'' (Ro 13:14). ``¿Quién acusará a los escogidos de Dios? Dios es el que justifica'' (Ro 8:33).
\end{column}
\begin{column}{0.49\textwidth}
\textbf{\color{azul}5. Amar su ley y hacer justicia}\par\small
``En la ley de \SENOR{} está su delicia'' (Sal 1:2); ``dulces más que miel'' (Sal 19:10).
Balanzas exactas (Lv 19:36); sin acepción de personas (Lv 19:15; Stg 2:1); defender al huérfano y a la viuda (Is 1:17). ``Bienaventurados los que tienen hambre y sed de justicia'' (Mt 5:6).
\end{column}
\end{columns}\par\bigskip
\begin{cita}Si sabéis que él es justo, sabed también que todo el que hace justicia es nacido de él.\fuente{(1 Jn 2:29, RVR1960)}\end{cita}
\end{frame}

% ------------------------------------------------------------------
\section{Síntesis}

\begin{frame}{Síntesis}
\scriptsize\renewcommand{\arraystretch}{1.18}
\begin{tabular}{@{}>{\RaggedRight\arraybackslash}p{1.6cm}>{\RaggedRight\arraybackslash}p{10.6cm}@{}}
\toprule
\textbf{Mirada} & \textbf{Aporte al atributo} \\ \midrule
Hebrea & Dios es justo (Sal 145:17; Dt 32:4); su justicia decantada en \emph{mishpatim}, \emph{edot}, \emph{juqqim} (Dt 4:45; 6:20; Sal 19; 119:142); Abraham los guardó (Gn 26:5); \emph{tsedeq} como norma; \emph{tsedaqá umishpat}; las ``justicias'' que salvan (Jue 5:11); la defensa del débil (Dt 10:18); el Juez de toda la tierra (Gn 18:25); el lamento que apela (Jer 12:1; Hab 2:4). \\
Judía & \emph{Ashré} diario (\emph{Berajot} 4b); ``Rey que ama la justicia'' (\emph{Berajot} 12b); ``si quieres un mundo\ldots'' (\emph{Bereshit Rabá} 39:6); \emph{tsidduq haddín} (\emph{Avodá Zará} 18a); el juez socio de Dios (\emph{Shabat} 10a); \emph{tsedaqá} = caridad; Qumrán (1QS XI). \\
Teológica & La justicia como límite que ningún atributo cruza (Is 5:16; Éx 34:7; Ro 5:21; 9:14; Pink, 2020, caps. 8, 13, 14, 19). \\
Paulina & El manto de justicia (Is 61:10) es Cristo (Gá 3:27); ley santa, justa y buena (Ro 7:12; 8:4); revelada (Ro 1:17); el justo y el que justifica (Ro 3:26); el intercambio (2 Co 5:21); Cristo hecho justicia (1 Co 1:30--31 con Jer 9:24); el justo juicio (Ro 2:5; Hch 17:31). \\
\bottomrule
\end{tabular}\par\medskip
\begin{cita}Mas alábese en esto el que se hubiere de alabar: en entenderme y conocerme, que yo soy \SENOR, que hago misericordia, juicio y justicia en la tierra\ldots\fuente{(Jer 9:24 [9:23], RVR1960)}\end{cita}
\end{frame}

\begin{frame}{Para la reflexión}
\small
\begin{enumerate}
\item ¿Qué injusticia que hoy no tiene solución humana necesito dejar en manos del ``Juez de toda la tierra'' (Gn 18:25; Ro 12:19)?
\item ¿Mi queja ante Dios es una apelación, como la de Jeremías, o una acusación, como la que Dios corrige en Job 40:8?
\item ¿Dónde estoy confiando todavía en ``mi propia justicia'' (Flp 3:9)?
\item ¿Qué ``balanza'' de mi trabajo, mi familia o mi iglesia necesita volver a ser exacta (Lv 19:36; Mi 6:8)?
\item ¿Quién es hoy, a mi alrededor, ``el huérfano y la viuda'' cuya causa Dios defiende (Dt 10:18; Jer 22:16)?
\end{enumerate}
\end{frame}

% ------------------------------------------------------------------
\section{Glosario}

\begin{frame}{Glosario}
\scriptsize
\begin{columns}[t,onlytextwidth]
\begin{column}{0.49\textwidth}
\begin{description}\setlength{\itemsep}{1pt}
\item[\emph{Dikaiokrisía}] (\grc{δικαιοκρισία}) ``justo juicio'' (Ro 2:5).
\item[\emph{Dikaiosýnē}] (\grc{δικαιοσύνη}) justicia (Ro 1:17).
\item[\emph{Dikaióō}] (\grc{δικαιόω}) justificar, declarar justo (Ro 3:26).
\item[\emph{Díkaios}] (\grc{δίκαιος}) justo; ``el Justo'' (Hch 3:14).
\item[\emph{Din}] (\heb{דִּין}) juicio; \emph{middat haddín}, la medida de la justicia.
\item[\emph{Hilastērion}] (\grc{ἱλαστήριον}) propiciación (Ro 3:25).
\item[\emph{Jasid}] (\heb{חָסִיד}) leal, fiel en amor (Sal 145:17).
\item[\emph{Mishpat}] (\heb{מִשְׁפָּט}) juicio, derecho (Gn 18:25).
\end{description}
\end{column}
\begin{column}{0.49\textwidth}
\begin{description}\setlength{\itemsep}{1pt}
\item[\emph{Shofet}] (\heb{שֹׁפֵט}) juez (Gn 18:25).
\item[\emph{Tsaddiq}] (\heb{צַדִּיק}) justo (Sal 145:17).
\item[\emph{Hitsdiq}] (\heb{הִצְדִּיק}) declarar justo (Éx 23:7; Is 53:11).
\item[\emph{Tsedaqá}] (\heb{צְדָקָה}) justicia; acto que salva; caridad.
\item[\emph{Tsedeq}] (\heb{צֶדֶק}) lo exacto, lo justo (Lv 19:36).
\item[\emph{Tsidduq haddín}] (\heb{צִדּוּק הַדִּין}) justificar el juicio de Dios.
\item[\emph{Tsidqot YHWH}] (\heb{צִדְקוֹת יְהוָה}) actos de salvación (Jue 5:11).
\item[\emph{YHWH tsidqenu}] (\heb{יְהוָה צִדְקֵנוּ}) ``el \SENOR, justicia nuestra'' (Jer 23:6).
\end{description}
\end{column}
\end{columns}
\end{frame}

% ------------------------------------------------------------------
\section{Referencias}

\begin{frame}{Referencias}
\fontsize{5.0}{6.0}\selectfont
\begin{columns}[t,onlytextwidth]
\begin{column}{0.49\textwidth}
\refapa{Bauer, W., Danker, F.~W., Arndt, W.~F., \& Gingrich, F.~W. (2000). \emph{A Greek-English lexicon of the New Testament and other early Christian literature} (3.ª ed.). University of Chicago Press.}
\refapa{Berkhof, L. (1996). \emph{Systematic theology} (Ed. combinada). Eerdmans. (Obra original publicada en 1938)}
\refapa{\emph{Bereshit Rabá} [8:5; 12:15; 39:6; 49:9]. (s.f.). Sefaria. https://www.sefaria.org}
\refapa{Brown, F., Driver, S.~R., \& Briggs, C.~A. (1906). \emph{A Hebrew and English lexicon of the Old Testament}. Clarendon Press.}
\refapa{Elliger, K., \& Rudolph, W. (Eds.). (1997). \emph{Biblia Hebraica Stuttgartensia} (5.ª ed.). Deutsche Bibelgesellschaft.}
\refapa{García Martínez, F. (1992). \emph{Textos de Qumrán}. Trotta.}
\refapa{Hays, R.~B. (1989). \emph{Echoes of Scripture in the letters of Paul}. Yale University Press.}
\refapa{Käsemann, E. (1969). The righteousness of God in Paul. En \emph{New Testament questions of today} (pp.~168--182). SCM Press.}
\refapa{Koehler, L., \& Baumgartner, W. (2001). \emph{The Hebrew and Aramaic lexicon of the Old Testament} (Ed. de estudio). Brill.}
\refapa{Lundbom, J.~R. (2004). \emph{Jeremiah 21--36} (Anchor Bible 21B). Doubleday.}
\refapa{Maimónides. (1904). \emph{The guide for the perplexed} (M. Friedländer, Trad.; 2.ª ed.). Routledge.}
\refapa{\emph{Misná: Berajot} [9:2, 9:5]. (s.f.). Sefaria. https://www.sefaria.org}
\refapa{Moo, D.~J. (1996). \emph{The epistle to the Romans}. Eerdmans.}
\refapa{Nestle, E., \& Aland, K. (Eds.). (2012). \emph{Novum Testamentum Graece} (28.ª ed.). Deutsche Bibelgesellschaft.}
\end{column}
\begin{column}{0.49\textwidth}
\refapa{Pink, A.~W. (2020). \emph{Los atributos de Dios} (E. Rodríguez Cruz, Trad.). Chapel Library. (Obra original publicada en 1930)}
\refapa{Qimron, E., \& Strugnell, J. (1994). \emph{Qumran Cave 4.V: Miqṣat maʿaśe ha-Torah} (DJD 10). Clarendon Press.}
\refapa{Rahlfs, A., \& Hanhart, R. (Eds.). (2006). \emph{Septuaginta} (Ed. rev.). Deutsche Bibelgesellschaft.}
\refapa{Rashi. (s.f.). \emph{Comentario a Génesis} [Gn 1:1]. Sefaria. https://www.sefaria.org}
\refapa{Sanders, J.~A. (1965). \emph{The Psalms Scroll of Qumrân Cave 11} (DJD 4). Clarendon Press.}
\refapa{\emph{Santa Biblia: Reina-Valera 1960}. (1960). Sociedades Bíblicas en América Latina.}
\refapa{\emph{Shemot Rabá} [3:6]. (s.f.). Sefaria. https://www.sefaria.org}
\refapa{\emph{Sidur Ashkenaz} [Amidá; \emph{Ashré}]. (s.f.). Sefaria. https://www.sefaria.org}
\refapa{Steinsaltz, A. (Ed.). (2017). \emph{Koren Talmud Bavli: Edición William Davidson} [Berajot 4b, 12b; Shabat 10a, 156b; Sanedrín 32b; Avodá Zará 18a; Bava Batra 75b]. Sefaria. https://www.sefaria.org}
\refapa{Strong, J. (1890). \emph{The exhaustive concordance of the Bible}. Hunt \& Eaton.}
\refapa{Vermes, G. (2004). \emph{The complete Dead Sea Scrolls in English} (Ed. rev.). Penguin.}
\refapa{Wenham, G.~J. (1994). \emph{Genesis 16--50} (WBC 2). Word.}
\refapa{Wright, N.~T. (2013). \emph{Paul and the faithfulness of God}. Fortress Press.}
\end{column}
\end{columns}
\par\smallskip{\color{gris}Las citas bíblicas son de la Reina-Valera 1960; ``Jehová'' se reemplaza por \SENOR. Hebreo: BHS; griego: NA28 y LXX de Rahlfs-Hanhart. Fuentes judías: Sefaria; lo rabínico se marca como \emph{precisión}.}
\end{frame}

\end{document}
